\documentclass[sigplan,twocolumn,review,anonymous]{acmart}

% EuroSys 2026 recommends the official ACM SIGPLAN/acmart proceedings setup.
% Replace the submission ID and conference metadata with the final values supplied by EuroSys/ACM.
\acmSubmissionID{PAPER-ID}
\setcopyright{none}
\settopmatter{printfolios=true,printacmref=false,printccs=false}
\renewcommand\footnotetextcopyrightpermission[1]{}
\pagestyle{plain}

\acmConference[EuroSys '26]{The European Conference on Computer Systems}{April 27--30, 2026}{Edinburgh, United Kingdom}
\acmYear{2026}
\copyrightyear{2026}
\acmDOI{}
\acmISBN{}

% acmart already loads compatible math fonts; avoid reloading amssymb under newtxmath.
\usepackage{amsmath}
\usepackage{algorithm}
\usepackage{algpseudocode}
\usepackage{mathtools}
\usepackage{adjustbox}
\usepackage{xspace}
\usepackage{balance}

\allowdisplaybreaks
\emergencystretch=2em
\sloppy

\makeatletter
\newenvironment{breakablealgorithm}
{%
  \par\addvspace{\baselineskip}%
  \begingroup\footnotesize
  \setlength{\baselineskip}{10pt}%
  \refstepcounter{algorithm}%
  \noindent\hrule\kern2pt%
  \renewcommand{\caption}[1]{%
    \noindent\textbf{Algorithm~\thealgorithm} ##1\par\kern2pt\hrule\kern2pt%
  }%
}
{%
  \kern2pt\hrule\par\endgroup\addvspace{\baselineskip}%
}
\makeatother
\algrenewcommand\algorithmicindent{0.7em}

\newcommand{\func}[1]{\operatorname{\mathsf{#1}}}
\newcommand{\hash}{\func{hash}}
\newcommand{\Cert}{\func{Cert}}
\newcommand{\NotarVote}{\func{NotarVote}}
\newcommand{\FirstVote}{\func{FirstVote}}
\newcommand{\FinalVote}{\func{FinalVote}}
\newcommand{\FinalSigners}{\func{FinalSigners}}
\newcommand{\NotarCert}{\func{NotarCert}}
\newcommand{\FirstCert}{\func{FirstCert}}
\newcommand{\FinalCert}{\func{FinalCert}}
\newcommand{\TimeoutCert}{\func{TimeoutCert}}
\newcommand{\ConflictProof}{\func{ConflictProof}}
\newcommand{\manyVotes}{\func{manyVotes}}
\newcommand{\allVotes}{\func{allVotes}}
\newcommand{\maxVotes}{\func{maxVotes}}
\newcommand{\id}{\mathsf{id}}
\newcommand{\Report}{\func{Report}}
\newcommand{\Slot}{\func{Slot}}
\newcommand{\Close}{\func{Close}}
\newcommand{\CloseCut}{\func{CloseCut}}
\newcommand{\CloseRecord}{\func{CloseRecord}}
\newcommand{\FreshCloseHash}{\func{FreshCloseHash}}
\newcommand{\CloseDead}{\func{CloseDead}}
\newcommand{\LiveCarry}{\func{LiveCarry}}
\newcommand{\CloseConflictProof}{\func{CloseConflictProof}}
\newcommand{\closeDecision}{\mathsf{closeDecision}}
\newcommand{\closeRound}{\mathsf{closeRound}}
\newcommand{\closeFinished}{\func{closeFinished}}
\newcommand{\started}{\func{started}}
\newcommand{\firstVoted}{\func{firstVoted}}
\newcommand{\finalVoted}{\func{finalVoted}}
\newcommand{\finished}{\func{finished}}
\newcommand{\enabled}{\func{enabled}}
\newcommand{\admissible}{\func{admissible}}
\newcommand{\admissibleFresh}{\func{admissibleFresh}}
\newcommand{\admissibleHistorical}{\func{admissibleHistorical}}
\newcommand{\dom}{\func{dom}}
\newcommand{\UpdateVisible}{\func{UpdateVisible}}
\newcommand{\StartClosing}{\func{StartClosing}}
\newcommand{\FinishClosing}{\func{FinishClosing}}
\newcommand{\valueOf}{\func{value}}
\newcommand{\certOutcome}{\func{certOutcome}}
\newcommand{\Terminal}{\func{Terminal}}
\newcommand{\ExtendsTree}{\func{ExtendsTree}}
\newcommand{\CompleteBlock}{\func{CompleteBlock}}
\newcommand{\ParentComplete}{\func{ParentComplete}}
\newcommand{\Complete}{\func{Complete}}
\newcommand{\timeoutVal}{\mathsf{timeout}}
\newcommand{\pending}{\mathsf{pending}}
\newcommand{\timeoutOutcome}{\mathsf{timeout}}
\newcommand{\notarized}{\func{notarized}}
\newcommand{\converged}{\func{converged}}
\newcommand{\fast}{\func{fast}}
\newcommand{\final}{\func{final}}
\newcommand{\notarKind}{\mathsf{notar}}
\newcommand{\fastKind}{\mathsf{fast}}
\newcommand{\finalKind}{\mathsf{final}}
\newcommand{\true}{\mathsf{true}}
\newcommand{\false}{\mathsf{false}}
\newcommand{\completeTree}{\mathsf{completeTree}}
\newcommand{\blockMap}{\mathsf{block}}
\newcommand{\accountOf}{\func{account}}
\newcommand{\frontier}{\mathsf{frontier}}
\newcommand{\ledgerRoot}{\mathsf{ledgerRoot}}
\newcommand{\closeHash}{\mathsf{closeHash}}
\newcommand{\carry}{\mathsf{carry}}
\newcommand{\Frontier}{\func{Frontier}}
\newcommand{\EpochBlocks}{\func{EpochBlocks}}
\newcommand{\parent}{\func{parent}}
\newcommand{\PayloadAvailable}{\func{PayloadAvailable}}
\newcommand{\ValidBlock}{\func{ValidBlock}}
\newcommand{\slotOf}{\func{slot}}
\newcommand{\finalizedSlot}{\mathsf{finalizedSlot}}
\newcommand{\FinalizeChain}{\func{FinalizeChain}}
\newcommand{\committee}{\mathsf{committee}}
\newcommand{\committeeAt}{\func{committeeAt}}
\newcommand{\Committees}{\func{Committees}}
\newcommand{\deriveCommittee}{\func{DeriveCommittee}}
\newcommand{\VotesIn}{\func{VotesIn}}
\newcommand{\NotarVotesIn}{\func{NotarVotesIn}}
\newcommand{\NotarizedBy}{\func{NotarizedBy}}
\newcommand{\timeoutReady}{\func{timeoutReady}}
\newcommand{\timeoutAllowed}{\func{timeoutAllowed}}
\newcommand{\FastDead}{\func{FastDead}}
\newcommand{\FinalDead}{\func{FinalDead}}
\newcommand{\TornFor}{\func{TornFor}}
\newcommand{\AllValuesDead}{\func{AllValuesDead}}
\newcommand{\Values}{\mathcal{V}}
\newcommand{\Reportable}{\func{Reportable}}
\newcommand{\reportStore}{\mathsf{reportStore}}
\newcommand{\ReportEquivocated}{\func{ReportEquivocated}}
\newcommand{\EligibleReport}{\func{EligibleReport}}
\newcommand{\votedFor}{\func{votedFor}}
\newcommand{\released}{\func{released}}
\newcommand{\safeToVote}{\func{safeToVote}}
\newcommand{\activeVote}{\func{activeVote}}
\newcommand{\closeCut}{\mathsf{cut}}
\newcommand{\closeState}{\mathsf{closeState}}
\newcommand{\governingHash}{\func{governingHash}}
\newcommand{\requiredFrontier}{\func{requiredFrontier}}
\newcommand{\DescendsFrom}{\func{DescendsFrom}}
\newcommand{\safeFromClose}{\func{safeFromClose}}
\newcommand{\StablePrefix}{\func{StablePrefix}}
\newcommand{\JointReportProof}{\func{JointReportProof}}
\newcommand{\ReportVisible}{\func{ReportVisible}}
\newcommand{\SelectedState}{\func{SelectedState}}
\newcommand{\FinalizedState}{\func{FinalizedState}}
\newcommand{\ReleasedState}{\func{ReleasedState}}
\newcommand{\cutVersion}{\mathsf{cutVersion}}
\newcommand{\frontierVersion}{\mathsf{frontierVersion}}
\newcommand{\ReconReq}{\func{ReconReq}}
\newcommand{\ReconDelta}{\func{ReconDelta}}
\newcommand{\ReconMiss}{\func{ReconMiss}}
\newcommand{\ApplyDelta}{\func{ApplyDelta}}
\newcommand{\seqOf}{\func{sequence}}
\newcommand{\ownerKey}{\mathsf{ownerKey}}
\newcommand{\signatureOf}{\func{signature}}
\newcommand{\sendsOf}{\func{sends}}
\newcommand{\receivesOf}{\func{receives}}
\newcommand{\representativeOf}{\func{representative}}
\newcommand{\bal}{\mathsf{bal}}
\newcommand{\repOf}{\mathsf{rep}}
\newcommand{\pendingSend}{\mathsf{pendingSend}}
\newcommand{\receivedSend}{\mathsf{receivedSend}}
\newcommand{\stake}{\func{stake}}
\newcommand{\weight}{\func{weight}}
\newcommand{\SignerWeight}{\func{SignerWeight}}

\title{RAI Block-Lattice Protocol: Minimal Specification}
\author{Anonymous Authors}
\affiliation{%
  \institution{Paper under review}
  \city{}
  \country{}
}
\email{}

\begin{abstract}
RAI applies one election mechanism to account-qualified slot, close-cut, and close-record elections. This specification defines a block-lattice ledger, owner-authorized account blocks, send and receive transitions, representative delegation, stake-weighted committees, epoch closing, safety, and liveness.
\end{abstract}

\keywords{block lattice, replication, consensus, stake-weighted elections, finality, Byzantine fault tolerance}

\begin{document}
\maketitle

\section{Overview}

RAI uses three election types.  A slot election \(\Slot(a,s,e)\) selects block evidence for account slot \((a,s)\)
in epoch \(e\).  A close-cut election \(\CloseCut(e,r)\) certifies a canonical set of visible
slot elections that epoch \(e\) must drain.  A close-record election \(\CloseRecord(e,r)\)
certifies the canonical finalized account frontiers through epoch \(e\), chained to the preceding
close-record hash.

Only three durable authenticated object classes are disseminated: blocks, signed reports, and signed
votes.  Certificates, report quorums, conflicts, close-round death, close cuts, frontier maps, close
records, outcomes, decisions, and installed close state are deterministic local facts derived from the
validated object store.  Correct replicas use duplicate-suppressing reliable gossip so that every valid
block, report, and vote learned by one correct online replica eventually becomes available to every
correct online replica.

Close votes carry only a canonical hash.  A replica retains every locally derived close-cut set and
frontier map for which it votes, indexed by that hash.  If another replica knows only a different
version, it reconstructs the target by requesting an add/remove delta from a replica that retains both
preimages, applies the delta, recomputes the target hash, and validates the reconstructed contents from
its locally available blocks, reports, and votes.  Reconciliation messages are transport metadata, not
consensus evidence.

A notarization never finalizes a slot by itself.  A slot becomes finalized through an ordinary global
fast/final certificate or through inclusion on an account chain committed by a decided close-record
hash.  Timeouts are local to one election and never skip a slot.  A same-account-slot retry is permitted only
after certified release.  Correct replicas vote for a descendant block only when every strict
non-genesis ancestor is already finalized, so released or merely complete blocks cannot later be
resurrected by ancestor finality.  Elections at distance at least two epochs are protected because every
vote in epoch \(e\) commits to and obeys the certified close state and finalized account frontiers of
epoch \(e-2\).  Safety further requires election-scoped votes, one first vote per correct replica and
committee, parent-first completion, a certified-stable ancestor prefix, one finalized block per account slot,
and certified release before a correct replica votes for a conflicting same-account-slot block.

\section{Objects, State, and Committees}

An account slot is a pair \(\lambda=(a,s)\), where \(a\in\mathsf{Accounts}\) and
\(s\ge1\) is the account-local sequence number.  An election id is either
\begin{center}
\begin{adjustbox}{max width=\columnwidth}
\(\displaystyle
    \id=\Slot(a,s,e),\qquad \id=\CloseCut(e,r),\qquad\text{or}\qquad
    \id=\CloseRecord(e,r).
\)
\end{adjustbox}
\end{center}
A slot-election non-timeout value is a block id \(b=\hash(B)\), where
\(\slotOf(B)=(a,s)\).  Every account \(a\) has an owner verification key
\(\ownerKey[a]\), with corresponding owner signing key \(sk_a\), fixed by genesis or by an
authenticated account-creation rule.  Only a signature verifiable under \(\ownerKey[a]\) can
authorize a non-genesis block on account \(a\).

A block has canonical body
\[
\operatorname{body}(B)=
\bigl(\slotOf(B),\parent(B),\sendsOf(B),\receivesOf(B),\representativeOf(B)\bigr).
\]
The send list contains entries \((d,x)\), where \(d\) is a destination account and \(x\) is a positive
integer amount.  A send is uniquely referenced by \(u=(\hash(B_s),i)\), where \(i\) is its canonical
index in \(\sendsOf(B_s)\).  The receive list contains such references.  The representative field is a
registered replica id.  The owner signature is
\[
\signatureOf(B)=\operatorname{Sign}_{sk_a}
\bigl(\hash(\mathsf{RAI/AccountBlock/v1}\parallel\operatorname{canon}(\operatorname{body}(B)))\bigr),
\]
and \(\hash(B)\) commits to both the canonical body and this signature.

The ledger state derived from a finalized block set is
\[
L=(\bal,\repOf,\pendingSend,\receivedSend).
\]
Here \(\bal[a]\) is the spendable balance of account \(a\), \(\repOf[a]\) is its latest
representative, \(\pendingSend[u]=(d,x)\) records an unreceived send, and
\(\receivedSend[u]\) records a consumed send reference.  This state is derived by replay from genesis;
it is not an independently trusted object.

For epoch \(e\), a close-cut version is a canonically ordered set
\[
C\subseteq\{\Slot(a,s,e):a\in\mathsf{Accounts},\ s\ge1\}.
\]
Its close-cut value is
\[
h_C=\hash\bigl(\mathsf{RAI/CloseCut}\parallel e\parallel\operatorname{canon}(C)\bigr).
\]
The voted value is only \(h_C\); the set \(C\) is local derived state retained as
\(\cutVersion[e,h_C]=C\).

For epoch \(e\), a frontier version is the canonically ordered map
\[
\Phi=\operatorname{sort}\{(a,f(a)):a\in\mathsf{Accounts}\},
\]
where every \(f(a)\) is a block hash.  The close-record value is
\[
h_F=\hash\bigl(
\mathsf{RAI/CloseRecord}\parallel e\parallel\closeHash[e-1]
\parallel\operatorname{canon}(\Phi)
\bigr).
\]
The voted value is only \(h_F\); the map \(\Phi\) is local derived state retained as
\(\frontierVersion[e,h_F]=\Phi\).  The local ledger-state root is the convenience value
\[
R^{\mathsf{ledger}}_e=\hash(\operatorname{canon}(\Phi_e)).
\]
Because every non-genesis block contains the hash of its predecessor, each frontier transitively
commits to every block on its account chain from genesis through that frontier, including blocks
finalized in previous epochs.  Replay of those chains deterministically commits to balances, pending
and received sends, and representatives.  The preceding close-record hash makes the decided close
records a hash chain over the complete ledger history.

\paragraph{Disseminated and derived objects.}
Blocks, reports, and votes are immutable authenticated protocol objects.  A correct replica advertises
and transfers each unique valid such object only to peers that do not already possess it.  A certificate
is not a wire object: it is the local fact that the validated vote store contains a quorum.  Joint report
proofs, conflicts, close-round death facts, close versions, outcomes, decisions, frontiers, release
state, and committees are likewise local derivations.  Inventory advertisements, retrieval requests,
and reconciliation deltas may be exchanged, but they carry no consensus authority.

\paragraph{Validated-close abstraction.}
A correct replica votes for a close hash only after it has locally constructed or reconstructed the
corresponding canonical set or frontier map and validated it using the underlying blocks, reports, and
votes.  Fresh admissibility is evaluated against the replica's current validated evidence before a new
close vote.  Once a close-cut hash obtains persistent same-hash certificate support in every applicable committee,
its hash-checked canonical preimage is historically admissible and never loses validity;
later report equivocation or newly learned evidence affects only fresh construction.  A historical
close-record version is validated against the immutable close input fixed by the decided cut, so later
ordinary finalizations outside that input cannot invalidate it.  Authority depends only on locally
derived vote quorums and the hash-checked canonical preimage, not on certificate or transport-message
arrival order.

After the close-cut decision, the authoritative cut is the version reconstructed for the decided hash.
The cut is supporting close-protocol state: it freezes the elections that may participate in draining
and frontier construction, but it is not the durable ledger-consensus value and is not an additional
field of the close-record hash.  After the close-record decision, the durable consensus checkpoint is
\(\closeHash[e]=h_F\), while \(\closeState[e]\) is reconstructed from the decided frontier map, the
decided cut, and the underlying validated blocks, reports, and votes.  Only this close result, not local
drain observations, feeds committee derivation and future slot-voting guards.  Each election id has its own timeout value
\(\timeoutVal_\id\).  In particular, \(\timeoutVal_{\Slot(a,s,e)}\) does not create an empty block,
finalize, or skip account slot \((a,s)\).  A later retry \(\Slot(a,s,e')\) is permitted only after the certified
close state releases the earlier election; a block finalized by the close record is not retried.

A replica stores known block data as
\begin{center}
\begin{adjustbox}{max width=\columnwidth}
\(\displaystyle
    \blockMap[\Slot(a,s,e),b]=B,
\)
\end{adjustbox}
\end{center}
when it accepts or learns valid data \(B\) for \(b\) in that slot election.
Uninitialized map entries have the following defaults:
\begin{center}
\begin{adjustbox}{max width=\columnwidth}
\(\displaystyle
\begin{gathered}
    epochStart[e]=\bot,
    \quad T_{\mathrm{election}}[\id]=\bot,
    \quad outcome[\id]=\pending,\\
    \closeDecision[I]=\bot,
    \qquad \closeRound[I]=0,
    \quad I\in\{\CloseCut(e),\CloseRecord(e)\},\\
    notarized[P_i,Q,\id]=secondLook[P_i,Q,\id]=\{\},
    \quad firstVote[P_i,\id,Q]=\bot,
    \quad finalVoted[P_i,\id,Q]=\false,\\
    \votedFor[P_i,\id,b]=\false,
    \quad \blockMap[\cdot,\cdot]=\bot,
    \quad \committee[\cdot]=\bot,
    \quad \closeHash[e]=\bot,
    \quad \ledgerRoot[e]=\bot,
    \quad \closeState[e]=\bot,
    \quad \cutVersion[e]=\{\},
    \quad \frontierVersion[e]=\{\},
    \quad \carry[\id]=\bot,
    \quad \closeCut[e]=\bot.
\end{gathered}
\)
\end{adjustbox}
\end{center}
Account-slot finality is tracked by
\begin{center}
\begin{adjustbox}{max width=\columnwidth}
\(\displaystyle
    \finalizedSlot[a,s]\in\{\bot\}\cup
    \{b:\exists B,\ b=\hash(B),\ \slotOf(B)=(a,s)\},
\)
\end{adjustbox}
\end{center}
initially \(\finalizedSlot[\cdot,\cdot]=\bot\).  A block is finalized exactly when its account slot
maps to its hash in \(\finalizedSlot\); no separate block-finality map is maintained.

Derived predicates and outcomes are
\begin{center}
\begin{adjustbox}{max width=\columnwidth}
\(\displaystyle
    \started[\id]\iff T_{\mathrm{election}}[\id]\neq\bot,
    \qquad
    \firstVoted(P_j,\id,Q)\iff firstVote[P_j,\id,Q]\neq\bot,
    \qquad
    \finalVoted(P_j,\id,Q)\iff finalVoted[P_j,\id,Q]=\true,
    \qquad
    \finished(\id)\iff outcome[\id]\neq\pending,
\)
\end{adjustbox}
\end{center}
\begin{center}
\begin{adjustbox}{max width=\columnwidth}
\(\displaystyle
    outcome[\id]\in
    \{\pending,\ \notarized(x),\ \converged(x),\ \timeoutOutcome,\ \fast(x),\ \final(x)\}.
\)
\end{adjustbox}
\end{center}
For a close instance \(I\), the variable
\[
    \closeDecision[I]\in
    \{\bot\}\cup\{(r,h)\}
\]
is instance-wide rather than round-local.  Here \(r\) is the deciding close round and \(h\) is the
decided hash.  The deciding certificate is derived locally from signed votes, and the canonical
preimage is retained separately in \(\cutVersion\) or \(\frontierVersion\).

The round-local variable \(outcome[I_r]\) records only local progress in round \(r\).  It is not the
authoritative close decision and does not prevent a replica from processing a valid fast or final
certificate learned later as a decision.  Define
\[
    \closeFinished(I)\iff\closeDecision[I]\neq\bot.
\]
Local completion is represented by membership in \(\completeTree\), not by an election
outcome.  A complete block is neither finished nor settled by completion alone.  A non-timeout slot
notarization certificate may finish a slot election as \(\notarized(b)\) for close drain, but
\(\notarized(b)\) does not finalize the slot.  At close, each replica uses its currently known valid votes and blocks to derive a fresh canonical
frontier map.  Global fast or final evidence is ordinary finality; locally derived timeout or conflicting
notarizations propose omission and release; and exactly one currently known non-timeout notarization,
with no known final or release evidence, may propose \(b\) for inclusion on a candidate account chain
ending at the proposed frontier.  Candidate inclusion is not finality.  Block \(b\) finalizes
only when the close-record hash committing to that frontier map is stored in
\(\closeDecision[\CloseRecord(e)]\) with a valid fast or final certificate.  A different decided frontier
or a decided omission releases \(b\), and later outside evidence cannot reverse the write-once decision.

A \(\FirstCert(\Slot(a,s,e),b)\), \(\NotarCert(\Slot(a,s,e),b)\), and
\(\FinalCert(\Slot(a,s,e),b)\) are compatible certificates for the same block \(b\).  A fast or final
certificate for \(b\) is incompatible with a notarization certificate for any \(b'\neq b\), including
the timeout value.  Conversely, two conflicting notarization certificates prove that no value can now
or later obtain an ordinary global fast or final certificate for that election; before close
certification they support release, while after a selecting close certificate they are losing
non-final evidence.  If a replica records an election outcome for the first time, its replica-local
\(outcome[\id]\) changes once from \(\pending\) to a non-pending value.  This variable is not an
agreed consensus value: correct replicas may temporarily record different terminal observations for
the same election, and no safety-critical retry or finality decision is derived from it without the
corresponding verifiable certificates and certified close state.  Later fast or final slot certificates
can still finalize their block before the close record decides the slot.  For
\(y\in\{\notarized(x),\fast(x),\final(x)\}\), \(\valueOf(y)=x\).  Election vote state is
logical committee-local state: \(firstVote[P_i,\id,Q]\), \(finalVoted[P_i,\id,Q]\),
\(notarized[P_i,Q,\id]\), and \(secondLook[P_i,Q,\id]\) are maintained independently for every
\(Q\in\Committees(\id)\).  Define the complete signed support of \(P_i\) in that committee instance by
\[
\func{Support}(P_i,Q,\id)=
notarized[P_i,Q,\id]\cup\bigl(\{firstVote[P_i,\id,Q]\}\setminus\{\bot\}\bigr).
\]
The first vote, including a timeout first vote, is therefore part of the support checked by every
final-vote rule.  A correct replica may final-vote \(x\) only when
\(\func{Support}(P_i,Q,\id)\subseteq\{x\}\).  An unscoped wire vote is only a compression of identical valid scoped
votes for every applicable committee in which the signer is a member; upon sending or receiving it,
a replica applies the corresponding state transition separately to each such committee.  Thus no
vote or lock in one committee implicitly changes vote state in another committee.

Let \(G\) be the genesis epoch state.  Replay of \(G\) supplies genesis balances and
representatives.  For any certified ledger state through epoch \(k\), define delegated stake
\[
\stake_k(P_i)=
\sum_{\substack{a\in\mathsf{Accounts}\\\repOf_k[a]=P_i}}\bal_k[a].
\]
The exact-proportional committee snapshot is the weighted replica map
\[
Q_k=\{(P_i,\stake_k(P_i)):\stake_k(P_i)>0\},
\qquad
\weight_{Q_k}(P_i)=\stake_k(P_i).
\]
Thus every positive-stake representative is a committee member and its consensus weight is exactly the
total finalized balance delegated to it.  For any signer set \(S\subseteq Q\), write
\[
\SignerWeight_Q(S)=\sum_{P_i\in S}\weight_Q(P_i),
\qquad W_Q=\SignerWeight_Q(Q).
\]
Committee membership and weights are frozen for the lifetime of the snapshot.  Let
\begin{center}
\begin{adjustbox}{max width=\columnwidth}
\(\displaystyle
    \committee_{\mathsf{gen}}:=\deriveCommittee(\mathsf{genesis},G)=Q_{\mathsf{gen}}.
\)
\end{adjustbox}
\end{center}
Each closed epoch derives
\begin{center}
\begin{adjustbox}{max width=\columnwidth}
\(\displaystyle
    \committee[e]:=\deriveCommittee(e,\closeHash[e],\closeState[e])=Q_e.
\)
\end{adjustbox}
\end{center}
The close hash identifies the certified snapshot, while replay of \(\closeState[e]\) supplies the
balances and representatives from which the weight map is deterministically reconstructed.
The certified close state that participates in every election of epoch \(e\) is identified by
\begin{center}
\begin{adjustbox}{max width=\columnwidth}
\(\displaystyle
\governingHash(e)=
\begin{cases}
\hash(G), & e<2,\\
\closeHash[e-2], & e\ge 2.
\end{cases}
\)
\end{adjustbox}
\end{center}
For \(e\ge2\), a replica may validate or emit an epoch-\(e\) vote only after it has obtained and
validated \(\closeState[e-2]\) and its close certificate.  This governing context is implicit: it is
not a serialized vote field and is not separately signed.  A verifier derives
\(\governingHash(e)\) from its certified close history and evaluates the vote, close-version data, and
committee lookup under that context; if the required certified close state is unavailable or
inconsistent, the vote is invalid for that verifier.
Committee lookup is
\begin{center}
\begin{adjustbox}{max width=\columnwidth}
\(\displaystyle
\committeeAt(k)=
\begin{cases}
\committee_{\mathsf{gen}}, & k<0,\\
\committee[k], & k\ge 0 \text{ and } \committee[k]\neq\bot,\\
\bot, & k\ge 0 \text{ and } \committee[k]=\bot.
\end{cases}
\)
\end{adjustbox}
\end{center}
Election committees are
\begin{center}
\begin{adjustbox}{max width=\columnwidth}
\(\displaystyle
    \Committees(\Slot(a,s,e))=\{\committeeAt(e-3),\committeeAt(e-2)\},
    \qquad
    \Committees(\CloseCut(e,r))=\{\committeeAt(e-3),\committeeAt(e-2)\},
    \qquad
    \Committees(\CloseRecord(e,r))=\{\committeeAt(e-3),\committeeAt(e-2)\},
\)
\end{adjustbox}
\end{center}
with duplicate committees collapsed.  Each \(Q\in\Committees(\id)\) runs the election core
independently.  A vote for \(\id\) is valid only from a replica in
\(\bigcup_{Q\in\Committees(\id)}Q\), and only if every committee lookup used by
\(\Committees(\id)\) is non-\(\bot\).  An unscoped vote for \(\id\) applies to every committee in
which the signer is a member; a committee-scoped vote \((\id,Q)\) applies only to \(Q\).  A replica
uses an unscoped vote exactly when the same action is enabled in all of its election committees, and
uses a scoped vote only when their enabled actions differ.  An unscoped vote occupies the
corresponding vote slot, and any post-final lock, in every committee where it applies; scoped and
unscoped copies from one signer are deduplicated per committee.  Every signed vote commits to
\(\id\) and either the complete committee set or the explicit committee \(Q\); its governing close
state is the implicit validation context derived from the vote's epoch.  This lagged committee
lookup permits epoch \(e+1\) to open as soon as epoch \(e\) enters
closing: its governing close state is epoch \(e-1\), which is already closed by the precondition of
\(\StartClosing(e)\), while \(\committee[e]\) is derived only when epoch \(e\) later finishes
closing.

\section{Voting Guards and Release}

A replica may cast first votes, timeout first votes, extra notarization votes, and timeout
notarization votes only while the election is enabled:
\begin{center}
\begin{adjustbox}{max width=\columnwidth}
\(\displaystyle
\begin{aligned}
\enabled(\Slot(a,s,e))\iff{}& \finalizedSlot[a,s]=\bot\\
&\wedge\bigl(state[e]=\mathsf{open}
\vee (state[e]=\mathsf{closing}\wedge\closeCut[e]\neq\bot\wedge\Slot(a,s,e)\in\closeCut[e])\bigr),\\
\enabled(\CloseCut(e,r))\iff{}& state[e]=\mathsf{closing},\\
\enabled(\CloseRecord(e,r))\iff{}& state[e]=\mathsf{closing}.
\end{aligned}
\)
\end{adjustbox}
\end{center}
Thus a slot election is normally vote-enabled only while its epoch is open, but after a certified
close cut is installed it remains vote-enabled during close-drain exactly when it is included in that
cut.

For a correct replica \(P_j\), casting any non-timeout first, notarization, or final vote for block
\(b\) in \(\Slot(a,s,e)\) records a unreleased vote:
\begin{center}
\begin{adjustbox}{max width=\columnwidth}
\(\displaystyle
    \votedFor[P_j,\Slot(a,s,e),b]=\true.
\)
\end{adjustbox}
\end{center}
An earlier unreleased vote is released only by a certified close decision:
\begin{center}
\begin{adjustbox}{max width=\columnwidth}
\(\displaystyle
\begin{aligned}
\released_{P_j}(q,b)
\iff{}& q=\Slot(a,s,e)\wedge state[e]=\mathsf{closed}\wedge\closeCut[e]\neq\bot\\
&{}\wedge\Bigl(q\notin\closeCut[e]\\
&\qquad\vee\ \neg\func{LedgerIncludes}_e(q,b)\Bigr).
\end{aligned}
\)
\end{adjustbox}
\end{center}
Define \(\func{LedgerIncludes}_e(q,b)\) to hold when \(q=\Slot(a,s,e)\),
\(\blockMap[q,b]=B\), and \(B\) occurs at logical account slot \((a,s)\) on the complete account chain from
genesis through the frontier committed by \(R^{\mathsf{ledger}}_e\).  For an included slot, timeout
and known-conflict evidence are positive witnesses permitting omission from the candidate ledger.  If
the certified ledger instead contains \(b'\), every other block \(b\neq b'\) for that slot is
released even when its notarization certificate is learned later.  If no candidate of \(q\) occurs on
the certified ledger, the election is released.  Raw timeout or conflict evidence may finish the slot
election for drain but does not by itself release future conflicting votes while the close record is
unsettled.  A local belief that an election is stale, idle, or disabled is not release.

The expanded certified close state is derived from the certified cut and ledger.  It maps an election
\(q\) to \(\FinalizedState(b,\pi_{\mathsf{ledger}})\) when
\(\func{LedgerIncludes}_e(q,b)\), and otherwise to \(\ReleasedState(\pi)\), with close-exclusion,
timeout, or conflict evidence derived from the validated report and vote store.  The state is cumulative with
respect to finality: every previously finalized block must occur on the certified account chains, and no
later close state may replace it with a conflicting value.  For account \(a\), define the governing
certified frontier
\begin{center}
\begin{adjustbox}{max width=\columnwidth}
\(\displaystyle
\requiredFrontier(a,e)=
\begin{cases}
\frontier[e-2,a], & e\ge2\text{ and }\closeState[e-2]\text{ is certified},\\
\bot, & \text{otherwise.}
\end{cases}
\)
\end{adjustbox}
\end{center}
Write \(\DescendsFrom(B,b_c)\) when \(b_c\) occurs on the strict parent chain of \(B\).  Define
\begin{center}
\begin{adjustbox}{max width=\columnwidth}
\(\displaystyle
\StablePrefix(B)\iff
\forall B_a\text{ on the strict non-genesis parent chain of }B:\
\finalizedSlot[\slotOf(B_a)]=\hash(B_a).
\)
\end{adjustbox}
\end{center}
Then
\begin{center}
\begin{adjustbox}{max width=\columnwidth}
\(\displaystyle
\begin{aligned}
\safeFromClose(\Slot(a,s,e),b)\iff{}& \exists B:\ b=\hash(B)\wedge\blockMap[\Slot(a,s,e),b]=B\\
&\wedge\StablePrefix(B)\\
&\wedge\bigl(e<2\ \vee\ (\closeState[e-2]\text{ is certified}
\wedge\DescendsFrom(B,\requiredFrontier(\accountOf(B),e)))\bigr).
\end{aligned}
\)
\end{adjustbox}
\end{center}
The \(\StablePrefix\) check is an account-chain vote lock.  It applies in every epoch, including the
adjacent open epoch while its predecessor is closing: a merely complete or unresolved ancestor cannot
support a valid descendant vote until it is finalized.  During close drain for epoch \(e\),
\(\safeFromClose\) is additionally evaluated against \(\mathsf{CloseInput}_e\): every ancestor beyond
\(\frontier[e-1,a]\) must be the selected block of an election in \(\closeCut[e]\), and cross-account
receive references must resolve in the corresponding close-local replay state.  Thus an epoch-\(e\)
drain vote cannot depend on a concurrently finalized epoch-\(e+1\) block.  A released ancestor never
satisfies this predicate.  A certified close record freezes every old election in its cut: correct
replicas emit no later vote in that closed election, and future votes for the account may extend only
the certified finalized frontier.

Earlier safety-relevant vote history for account slot \((a,s)\) is
\begin{center}
\begin{adjustbox}{max width=\columnwidth}
\(\displaystyle
\begin{aligned}
\activeVote_{P_j}(a,s,e_0,b)
\iff{}& \votedFor[P_j,\Slot(a,s,e_0),b]=\true.
\end{aligned}
\)
\end{adjustbox}
\end{center}
A later non-timeout slot vote is safe exactly when
\begin{center}
\begin{adjustbox}{max width=\columnwidth}
\(\displaystyle
\begin{aligned}
\safeToVote_{P_j}(\Slot(a,s,e'),b')
\iff{}& \finalizedSlot[a,s]\in\{\bot,b'\} \\
&\wedge\safeFromClose(\Slot(a,s,e'),b')\\
&\wedge\forall e_0<e',\ \forall b\neq b':
\activeVote_{P_j}(a,s,e_0,b)\Rightarrow
\released_{P_j}(\Slot(a,s,e_0),b).
\end{aligned}
\)
\end{adjustbox}
\end{center}
Thus a correct replica may vote again in the same election for the same block as allowed by the
election core, but may vote in a later same-account-slot election only after every earlier obligation has been
certifiably released.  A finalized slot is not retried.  For later slots of the same account,
\(\safeFromClose\) requires every valid candidate to descend from the governing certified finalized
frontier and to have only finalized strict ancestors.  Timeout and close votes do not create slot
unreleased votes.

\section{Certificates and Vote Counting}

RAI instantiates, for each fixed election identifier and weighted committee, the vote-counting core of
Kudzu~\cite{shoup2025kudzu}.  Kudzu supplies the single-committee counting pattern behind its
fast-path safety lemmas and its two-case termination argument.  RAI replaces voter cardinality with the
frozen consensus weight of the applicable stake snapshot.  It does not import Kudzu's block proposal,
reconstruction, or slot-advance semantics.  Election-scoped admissibility, multiple-committee
composition, close-round carry, close-version reconciliation, and cross-epoch reasoning are RAI-specific
and are stated and proved below.

\paragraph{Committee-local election kernel.}
Fix an election identifier \(\id\) and a weighted committee \(Q\in\Committees(\id)\).  Every vote
signature binds the vote kind, \(\id\), \(Q\), and value.  A correct member sends at most one first
vote in \((\id,Q)\), and a first vote for \(x\) also counts as a notarization vote for \(x\).  A
correct member may issue a second-look notarization vote for a non-timeout value \(x\) only after
observing first-vote signer weight strictly greater than \(F_Q+P_Q\) for \(x\) and validating
\(\admissible(\id,x)\).  It may issue a timeout notarization vote only when
\(\timeoutReady(Q,\id)\) holds.  Before final-voting \(x\), its complete committee-local support is a
subset of \(\{x\}\); after final-voting, it emits no later first, notarization, timeout, or final vote
in that committee instance.  Valid votes and certificates are persistent.  These rules are the complete
committee-local assumptions used when adapting Kudzu's Lemmas~5.4--5.7 and Lemma~5.10; no other
property of Kudzu is assumed.

Each committee \(Q\in\Committees(\id)\) has total weight \(W_Q\), maximum Byzantine weight
\(F_Q\), and quorum-slack weight \(P_Q\).  Assume
\[
P_Q\ge F_Q,
\qquad
W_Q>3F_Q+2P_Q.
\]
A local non-timeout certificate \(\Cert_Q(k,\id,x)\), where
\(k\in\{\notarKind,\fastKind,\finalKind\}\) and \(x\neq\timeoutVal_\id\), has threshold
\begin{center}
\begin{adjustbox}{max width=\columnwidth}
\(\displaystyle
\begin{array}{c|c|c}
\text{kind} & \text{signer-weight threshold in }Q & \text{value restriction}\\
\hline
\notarKind & \ge W_Q-F_Q-P_Q \text{ notarization weight} & \text{none}\\
\fastKind & \ge W_Q-P_Q \text{ first-vote weight} & x\neq\timeoutVal_\id\\
\finalKind & \ge W_Q-F_Q-P_Q \text{ final-vote weight} & x\neq\timeoutVal_\id.
\end{array}
\)
\end{adjustbox}
\end{center}
A certificate is the local fact that individually signed votes of sufficient distinct-signer weight are
present; it is not a separate BLS, threshold-signature, or extra vote-bearing object.  All votes counted
for a local certificate are valid for the same \((\id,Q,x)\), and each signer contributes its frozen
committee weight at most once.  A first vote for \((\id,Q,x)\) is also counted as a notarization vote
for that tuple.  Valid votes and certificates do not expire.

A local timeout certificate \(\TimeoutCert(\id,Q)\) exists when valid timeout notarization votes have
signer weight at least \(W_Q-F_Q-P_Q\).  The local terminal-result domain is
\[
L_Q(\id)\in\{\timeoutOutcome,\notarized(x),\fast(x),\final(x)\}.
\]
The RAI election result is the deterministic merge of the local results.  For the same non-timeout
value \(x\),
\begin{center}
\begin{adjustbox}{max width=\columnwidth}
\(\displaystyle
\begin{aligned}
\func{merge}(\notarized(x),\notarized(x))&=\notarized(x),\\
\func{merge}(\notarized(x),\fast(x))&=\func{merge}(\notarized(x),\final(x))=\notarized(x),\\
\func{merge}(\fast(x),\fast(x))&=\fast(x),\\
\func{merge}(\fast(x),\final(x))&=\func{merge}(\final(x),\final(x))=\final(x),\\
\func{merge}(\timeoutOutcome,y)&=\timeoutOutcome.
\end{aligned}
\)
\end{adjustbox}
\end{center}
The equations are symmetric.  Two non-timeout local results for different values, including two local
final results for different values, merge to \(\timeoutOutcome\).  Their certificate-level evidence is a
conflict proof that the round cannot produce a compatible global fast or final result.  With one
committee, the merged result is that committee's local result.  The global aliases
\(\NotarCert(\id,x)\), \(\FirstCert(\id,x)\), and \(\FinalCert(\id,x)\) mean that the merged result
is respectively \(\notarized(x)\), a joint fast result \(\fast(x)\) from every election committee, or
\(\final(x)\); \(\NotarCert(\id,\timeoutVal_\id)\) means that the merged result is
\(\timeoutOutcome\).

A merged result finishes an election round according to the slot/close terminal rules below.
Fast and final results are valid only for non-timeout values.  A matching non-timeout notarization
result in every close committee records round-local progress as \(\converged(x)\) and creates a live
carry unless positive death evidence is known.  A valid fast or final certificate in any close round
decides the logical close instance through \(\closeDecision\), regardless of the replica's current
round or the round-local value of \(outcome\).  Correct replicas continue to use the committee-local
election kernel above in close rounds.

\begin{center}
\begin{adjustbox}{max width=\columnwidth}
\(\displaystyle
\Terminal(t,\id,x)\iff
\begin{cases}
    t\in\{\notarKind,\fastKind,\finalKind\}, & \id=\Slot(a,s,e),\\
    t=\notarKind\vee(t\in\{\fastKind,\finalKind\}\wedge x\neq\timeoutVal_\id),
        & \id\in\{\CloseCut(e,r),\CloseRecord(e,r)\},
\end{cases}
\)
\end{adjustbox}
\end{center}
where \(t,x\) denote the kind and value represented by the merged result, and
\begin{center}
\begin{adjustbox}{max width=\columnwidth}
\(\displaystyle
\certOutcome(t,x)=
\begin{cases}
\timeoutOutcome, & t=\notarKind\wedge x=\timeoutVal_\id,\\
\notarized(x), & t=\notarKind\wedge x\neq\timeoutVal_\id\wedge\id=\Slot(a,s,e),\\
\converged(x), & t=\notarKind\wedge x\neq\timeoutVal_\id\wedge\id\in\{\CloseCut(e,r),\CloseRecord(e,r)\},\\
\fast(x), & t=\fastKind,\\
\final(x), & t=\finalKind.
\end{cases}
\)
\end{adjustbox}
\end{center}

For \(Q\in\Committees(\id)\), define
\begin{center}
\begin{adjustbox}{max width=\columnwidth}
\(\displaystyle
\VotesIn(Q,\id,x)=\{P_i\in Q: firstVote[P_i,\id,Q]=x\},
\qquad
\NotarVotesIn(Q,\id,x)=
\{P_i\in Q: firstVote[P_i,\id,Q]=x\vee\NotarVote(P_i,\id,Q,x)\in\mathsf{pool}\}.
\)
\end{adjustbox}
\end{center}
\begin{center}
\begin{adjustbox}{max width=\columnwidth}
\(\displaystyle
\allVotes(Q,\id)=
\SignerWeight_Q(\{P_i\in Q:firstVote[P_i,\id,Q]\neq\bot\}),
\)
\end{adjustbox}
\end{center}
\begin{center}
\begin{adjustbox}{max width=\columnwidth}
\(\displaystyle
\maxVotes(Q,\id)=\max_{x\neq\timeoutVal_\id}\SignerWeight_Q(\VotesIn(Q,\id,x)),
\)
\end{adjustbox}
\end{center}
where the maximum is \(0\) if there is no non-timeout value, and
\begin{center}
\begin{adjustbox}{max width=\columnwidth}
\(\displaystyle
\manyVotes(Q,\id)=
\{x\neq\timeoutVal_\id:\ \SignerWeight_Q(\VotesIn(Q,\id,x))>F_Q+P_Q\}.
\)
\end{adjustbox}
\end{center}
Timeout notarization for committee \(Q\) is enabled by
\begin{center}
\begin{adjustbox}{max width=\columnwidth}
\(\displaystyle
\timeoutReady(Q,\id)\iff
\allVotes(Q,\id)-\maxVotes(Q,\id)>F_Q+P_Q.
\)
\end{adjustbox}
\end{center}

The monotone slack rule above is the signed timeout trigger, following Kudzu's main loop
\cite{shoup2025kudzu}.  For diagnostics, let \(\Values(\id)\) be the finite election-local set of
non-timeout values that currently appear in valid close-version data or valid votes for \(\id\), and
define
\begin{center}
\begin{adjustbox}{max width=\columnwidth}
\(\displaystyle
\NotarizedBy(P_i,\id,Q)=
\{x:firstVote[P_i,\id,Q]=x\vee\NotarVote(P_i,\id,Q,x)\in\mathsf{pool}\}.
\)
\end{adjustbox}
\end{center}
\begin{center}
\begin{adjustbox}{max width=\columnwidth}
\(\displaystyle
\begin{aligned}
\FastDead(Q,\id,y)\iff{}&
\SignerWeight_Q(\{P_i\in Q:firstVote[P_i,\id,Q]\notin\{\bot,y\}\})>F_Q+P_Q,\\
\TornFor(Q,\id,y)={}&
\{P_i\in Q:\exists x\in\NotarizedBy(P_i,\id,Q),\ x\neq y\}
\setminus\{P_i\in Q:\FinalVote(P_i,\id,Q,y)\in\mathsf{pool}\},\\
\FinalDead(Q,\id,y)\iff{}&
\SignerWeight_Q(\TornFor(Q,\id,y))>2F_Q+P_Q,\\
\AllValuesDead(Q,\id)\iff{}&
\forall y\in\Values(\id):\FastDead(Q,\id,y)\wedge\FinalDead(Q,\id,y).
\end{aligned}
\)
\end{adjustbox}
\end{center}
\begin{center}
\begin{adjustbox}{max width=\columnwidth}
\(\displaystyle
\timeoutAllowed(Q,\id)\iff \timeoutReady(Q,\id).
\)
\end{adjustbox}
\end{center}
The set \(\Values(\id)\) is learned incrementally, so \(\AllValuesDead(Q,\id)\) is only a local
diagnostic over currently known values.  It does not authorize a signed timeout vote: an unknown value
could later acquire first-vote weight, and a timeout vote based on an incomplete candidate set would not
be covered by the weighted incompatibility argument.  A correct replica may use \(\AllValuesDead\) to
avoid futile local reconstruction after certificate-level death evidence is already present, but a
timeout notarization vote is issued only under the monotone \(\timeoutReady\) predicate.  Such a vote is
an ordinary signed notarization vote for \(\timeoutVal_\id\); receivers validate and count it like any
other notarization vote.

\paragraph{Lemma (committee-local certificate incompatibility).}
Fix \((\id,Q)\).  Assume Byzantine signers have total weight at most \(F_Q\) and
\(W_Q>3F_Q+2P_Q\).  If a valid local fast or final certificate exists for \(x\), then no valid local
notarization certificate exists for any \(y\neq x\), including \(\timeoutVal_\id\).

For the fast case, a fast certificate for \(x\) contains first-vote signer weight at least
\(W_Q-P_Q\).  Even allowing all Byzantine weight to equivocate, at most \(F_Q+P_Q\) distinct-signer
weight can be recorded outside \(x\).  Hence no different value can enter \(\manyVotes(Q,\id)\), whose
threshold is strictly greater than \(F_Q+P_Q\), and
\(\allVotes(Q,\id)-\maxVotes(Q,\id)\le F_Q+P_Q\), so \(\timeoutReady(Q,\id)\) is false.  Correct
members therefore send no second-look or timeout notarization vote for a value different from \(x\).
The remaining Byzantine weight is less than the notarization threshold \(W_Q-F_Q-P_Q\).

For the final case, a final certificate for \(x\) and a notarization certificate for \(y\neq x\) would
have signer weights at least \(W_Q-F_Q-P_Q\).  Their intersection has weight at least
\[
2(W_Q-F_Q-P_Q)-W_Q=W_Q-2F_Q-2P_Q>F_Q.
\]
It therefore contains correct signer weight, contradicting the support condition and post-final lock of
the committee-local kernel.  This is the stake-weighted RAI instantiation of Kudzu's fast-finalization
and incompatibility arguments~\cite[Lemmas~5.4--5.7]{shoup2025kudzu}.

Since a global fast or final certificate requires the same value in every election committee, any two
valid committee-local non-timeout terminal certificates for different values in applicable committees
form a conflict proof
\[
\ConflictProof(\id,x,y),\qquad x\neq y,
\]
where a terminal certificate may be notarization, fast, or final evidence.  Such evidence excludes every
possible present or future compatible global fast or final certificate for \(\id\).  A merged timeout
result records either an explicit timeout certificate or such incompatible terminal evidence; the
underlying signed votes remain in the validated object store as the release witness.

\paragraph{Lemma (committee composition).}
Suppose every \(Q\in\Committees(\id)\) has derived a local terminal result.  Because vote validity and
locks are interpreted per \((\id,Q)\), overlapping membership and unscoped wire compression do not
couple the logical committee instances.  The deterministic merge therefore produces a global terminal
result immediately.  It finalizes a non-timeout value only when every committee supplies a compatible
fast or final result for that same value.  Timeout in one committee, or different non-timeout values in
two committees, cannot produce finality.  By the preceding lemma, two different globally finalized
values are impossible.

\subsection{Close-round carry and death evidence}

Fix a close instance
\[
I\in\{\CloseCut(e),\CloseRecord(e)\},
\]
and write \(I_r\) for its election identifier in close round \(r\).

A close-round death proof is positive certificate evidence that round \(r\) cannot now or later
produce a deciding fast or final certificate:
\[
\CloseDead_{P_j}(I_r)\iff
\begin{cases}
\NotarCert(I_r,\timeoutVal_{I_r})\in\mathsf{pool}_{P_j},\\
\text{or }P_j\text{ has }\CloseConflictProof(I_r,h,h')\text{ for }h\neq h'.
\end{cases}
\]

A conflict proof consists of valid certificate-level committee-local terminal evidence, including
notarization, fast, or final evidence, whose coexistence is incompatible with every possible deciding
fast or final certificate for \(I_r\).  The underlying signed votes are retained and disseminated through ordinary vote gossip.  A local clock timeout, changed preference, or newly learned
uncertified evidence is not a death proof.

For a non-timeout hash \(h\), define
\[
\LiveCarry_{P_j}(I_r,h)\iff
\NotarCert(I_r,h)\in\mathsf{pool}_{P_j}
\wedge
\neg\CloseDead_{P_j}(I_r).
\]

The predicate is local.  Correct replicas may temporarily derive different live carries if a
the signed votes completing a conflicting quorum have not yet been completely gossiped.  Such disagreement is safe because the
hidden conflict makes the source round incapable of deciding.

\section{Blocks, Admissibility, and Finality}

The complete block tree \(\completeTree\) contains block data, with one account tree per account.  A
complete slot block is represented as \(B\), and \(\slotOf(B)=(a,s)\).  Define
\begin{center}
\begin{adjustbox}{max width=\columnwidth}
\(\displaystyle
\CompleteBlock(b)\iff
b\in\{\mathsf{genesis}_a:a\in\mathsf{Accounts}\}\vee\exists B\in\completeTree:b=\hash(B).
\)
\end{adjustbox}
\end{center}
\begin{center}
\begin{adjustbox}{max width=\columnwidth}
\(\displaystyle
\begin{aligned}
\ParentComplete(B,a,s)
\iff{}& (\parent(B)=\mathsf{genesis}_{a}\wedge s=1)\\
&\vee\exists B_p\in\completeTree:
\parent(B)=\hash(B_p)\wedge\slotOf(B_p)=(a,s-1).
\end{aligned}
\)
\end{adjustbox}
\end{center}
Here \(\slotOf(\mathsf{genesis}_a)=(a,0)\).  Define
\[
\func{AncestorsConsistent}(B)\iff
\forall B_a\text{ on the parent chain of }B:\
\finalizedSlot[\slotOf(B_a)]\in\{\bot,\hash(B_a)\}.
\]

A receive reference \(u=(\hash(B_s),i)\) is available only when \(B_s\) is finalized, index \(i\)
selects a send \((d,x)\in\sendsOf(B_s)\), and the finalized ledger state contains
\(\pendingSend[u]=(d,x)\).  Requiring the send block to be finalized prevents cyclic or same-close
cross-account funding dependencies.  For a candidate block \(B\) of account \(a\), let \(L_B\) be the
unique finalized ledger state against which its parent and receive references are evaluated.  When
\(B\) is validated for \(q=\Slot(a,s,e)\) during close drain, \(L_B\) is the close-local state rooted at
\(\closeState[e-1]\) and restricted by \(\mathsf{CloseInput}_e\); blocks, sends, and receives whose
elections are outside the decided cut are unavailable to that validation even if they have become
ordinary-finalized concurrently.  Define
\[
\operatorname{credit}(B,L_B)=
\sum_{u\in\receivesOf(B)}\pendingSend_{L_B}[u].\mathsf{amount},
\qquad
\operatorname{debit}(B)=\sum_{(d,x)\in\sendsOf(B)}x.
\]
The block-validity predicate is
\begin{center}
\begin{adjustbox}{max width=\columnwidth}
\(\displaystyle
\begin{aligned}
\ValidBlock(B)\iff{}& \slotOf(B)=(a,s)\wedge\ParentComplete(B,a,s)\\
&\wedge\operatorname{Verify}_{\ownerKey[a]}
 \bigl(\signatureOf(B),
 \hash(\mathsf{RAI/AccountBlock/v1}\parallel\operatorname{canon}(\operatorname{body}(B)))\bigr)\\
&\wedge\text{every send amount is a positive representable integer}\\
&\wedge\text{all receive references in }B\text{ are distinct}\\
&\wedge\forall u\in\receivesOf(B),\ \exists x>0:
 u\notin\receivedSend_{L_B}\wedge\pendingSend_{L_B}[u]=(a,x)\\
&\wedge\bal_{L_B}[a]+\operatorname{credit}(B,L_B)-\operatorname{debit}(B)\ge0\\
&\wedge\representativeOf(B)\in\mathsf{RegisteredReplicas}.
\end{aligned}
\)
\end{adjustbox}
\end{center}
Only the account owner can therefore produce an authorized block, an account balance can never become
negative, every send can be consumed at most once, and only its declared destination account can consume
it.  Applying a valid block produces
\[
\begin{aligned}
\bal'[a]&=\bal[a]+\operatorname{credit}(B,L_B)-\operatorname{debit}(B),\\
\repOf'[a]&=\representativeOf(B).
\end{aligned}
\]
Every received reference is removed from \(\pendingSend\) and added to \(\receivedSend\); every send
at index \(i\) creates
\[
\pendingSend'[(\hash(B),i)]=(d_i,x_i).
\]
With no mint, burn, or fee transition, replay also verifies the supply invariant
\[
\sum_a\bal[a]+\sum_{u\in\dom(\pendingSend)}\pendingSend[u].\mathsf{amount}
=\mathsf{GenesisSupply}.
\]
The replay result is independent of traversal order provided account-parent edges and finalized
send-before-receive edges are respected.

The complete block tree contains at most one entry for each block hash.  A slot value extends the
complete tree when
\begin{center}
\begin{adjustbox}{max width=\columnwidth}
\(\displaystyle
\begin{aligned}
\ExtendsTree(\Slot(a,s,e),b)
\iff{}& \exists B: b=\hash(B)\wedge\blockMap[\Slot(a,s,e),b]=B\\
&\wedge\slotOf(B)=(a,s)\wedge\ParentComplete(B,a,s)
\wedge\func{AncestorsConsistent}(B)\\
&\wedge\ValidBlock(B)\wedge\finalizedSlot[a,s]\in\{\bot,b\}.
\end{aligned}
\)
\end{adjustbox}
\end{center}
For a slot election \(q=\Slot(a,s,e)\), define a start witness by
\begin{center}
\begin{adjustbox}{max width=\columnwidth}
\(\displaystyle
\begin{aligned}
\func{StartWitness}_e(q)\iff{}&
\bigl(\exists Q\in\Committees(q):\mathsf{pool}\text{ contains a valid first vote for }q\text{ in }Q\bigr)\\
&\vee\bigl(\forall Q\in\Committees(q):
\SignerWeight_Q(\{P_i\in Q:q\in report[P_i,e]\})>F_Q\bigr).
\end{aligned}
\)
\end{adjustbox}
\end{center}
The witness is a local predicate over already gossiped reports and votes; it is not a separate object.

\paragraph{Stable admissibility and fresh preference.}
Admissibility is distinct from the value that a correct replica currently prefers to first-vote.  For
correct replica \(P_j\), let \(E_j\) be its validated blocks, reports, and votes.  Canonical fresh
construction must satisfy
\[
E_i=E_j
\Longrightarrow
\FreshCloseHash_{P_i}(I)=\FreshCloseHash_{P_j}(I).
\]
Serialization order, witness choice, reconciliation path, and other transport choices do not affect a
close hash.

Slot admissibility remains
\begin{center}
\begin{adjustbox}{max width=\columnwidth}
\(\displaystyle
\admissible(\Slot(a,s,e),b)\iff
b\neq\timeoutVal_{\Slot(a,s,e)}\wedge\ExtendsTree(\Slot(a,s,e),b)
\wedge\safeFromClose(\Slot(a,s,e),b).
\)
\end{adjustbox}
\end{center}

A reconstructed close-cut version \(C\) is freshly admissible for \(h_C\) when
\begin{center}
\begin{adjustbox}{max width=\columnwidth}
\(\displaystyle
\begin{aligned}
\admissibleFresh(\CloseCut(e),h_C,C)\iff{}&
 h_C=\hash(\mathsf{RAI/CloseCut}\parallel e\parallel\operatorname{canon}(C))\\
&\wedge\exists\rho:\JointReportProof(e,\rho)\\
&\wedge\{\Slot(a,s,e):(a,s)\in\ReportVisible(e,\rho)\}\subseteq C\\
&\wedge\forall q\in C:\func{StartWitness}_e(q).
\end{aligned}
\)
\end{adjustbox}
\end{center}
Define persistent certificate support for a cut hash by
\begin{center}
\begin{adjustbox}{max width=\columnwidth}
\(\displaystyle
\begin{aligned}
\func{CutSupported}(e,h_C)\iff{}&\exists r:\ \forall Q\in\Committees(\CloseCut(e,r)),\\
&\exists k_Q\in\{\notarKind,\fastKind,\finalKind\}:\
\Cert_Q(k_Q,\CloseCut(e,r),h_C).
\end{aligned}
\)
\end{adjustbox}
\end{center}
This predicate is monotone because the underlying committee-local certificates and their signed votes
do not expire, even if later incompatible evidence changes the currently merged round result.
Historical cut admissibility and the combined predicate are
\begin{center}
\begin{adjustbox}{max width=\columnwidth}
\(\displaystyle
\begin{aligned}
\admissibleHistorical(\CloseCut(e),h_C,C)\iff{}&
 h_C=\hash(\mathsf{RAI/CloseCut}\parallel e\parallel\operatorname{canon}(C))\\
&\wedge\func{CutSupported}(e,h_C),\\
\admissible(\CloseCut(e),h_C,C)\iff{}&
 \admissibleFresh(\CloseCut(e),h_C,C)\\
&\vee\admissibleHistorical(\CloseCut(e),h_C,C).
\end{aligned}
\)
\end{adjustbox}
\end{center}
The proof \(\rho\) and each start witness are selected locally from the report and vote store.  They are
not serialized with \(C\), and close votes continue to sign only the close hash rather than a witness
bundle reference.  Additional non-conflicting reports may change the replica's current visible set, and
newly learned report equivocation removes that signer's weight from fresh report calculations.  A
retained version is freshly revalidated before unsupported voting; once \(\func{CutSupported}(e,h_C)\)
holds, the historical branch is permanent and later report-store changes cannot invalidate the cut.

For a cut election \(q\), let \(\Sigma_{j,q}\) be all valid votes for \(q\) currently known to
\(P_j\).  Fresh close-record construction applies the deterministic resolver
\[
\operatorname{resolve}(q,\Sigma_{j,q})=
\begin{cases}
\FinalizedState(b,\pi),
  & \Sigma_{j,q}\text{ derives valid global fast or final evidence for }b,\\
\ReleasedState(\pi_{\mathsf{timeout}}),
  & \Sigma_{j,q}\text{ derives valid timeout evidence},\\
\ReleasedState(\pi_{\mathsf{conflict}}),
  & \Sigma_{j,q}\text{ derives conflicting notarization evidence},\\
\SelectedState(q,b,\pi_{\mathsf{notar}}),
  & b\text{ is the unique known non-timeout notarized value}\\
  & \text{and no final or release evidence is currently known},\\
\bot, & \text{otherwise.}
\end{cases}
\]
The symbols \(\pi\) denote locally selected signed-vote witnesses, not certificate objects.

Once the decided cut is installed, it fixes the immutable close input
\begin{center}
\begin{adjustbox}{max width=\columnwidth}
\(\displaystyle
\mathsf{CloseInput}_e=
\bigl(\closeHash[e-1],\{\frontier[e-1,a]:a\in\mathsf{Accounts}\},\closeCut[e]\bigr).
\)
\end{adjustbox}
\end{center}
This is an objective protocol boundary rather than a replica-local wall-clock snapshot.  Starting from
the preceding decided frontiers, the resolver is applied only to elections in \(\closeCut[e]\), and the
validated ancestry closure derives one candidate frontier for every account.  A reconstructed frontier
map \(\Phi\) is admissible for \(h_F\) when
\[
\begin{aligned}
\admissible(\CloseRecord(e),h_F,\Phi)\iff{}&
 h_F=\hash(\mathsf{RAI/CloseRecord}\parallel e\parallel\closeHash[e-1]
       \parallel\operatorname{canon}(\Phi))\\
&\wedge\func{validatedFrontiers}(e,\closeCut[e],\Phi).
\end{aligned}
\]
The predicate \(\func{validatedFrontiers}\) obtains every block on every account chain from genesis
through the claimed frontiers, validates every block and transition, preserves every block committed by
the preceding decided frontiers, and permits at most one block per logical account slot.  For every
\(q\in\closeCut[e]\), a block included on \(\Phi\) must have locally derivable global fast/final evidence
or a non-timeout notarization certificate for that same block; an omitted election must have locally
derivable timeout or conflicting terminal-certificate evidence.  Every strict ancestor of a selected
block must occur on the preceding certified account chain or be selected through an election in the
same frozen \(\closeCut[e]\).  Only elections in
the frozen \(\closeCut[e]\) may advance an account beyond \(\frontier[e-1,\cdot]\).  An ordinary-finalized
block whose election is outside \(\closeCut[e]\), including a block elected in the concurrently open
epoch \(e+1\), is neither required nor permitted to advance \(\Phi\) and first becomes eligible for a
later close record.  Additional evidence for elections in the frozen input may change the fresh
canonical frontier map, but it does not invalidate a historical map whose required positive witnesses
remain available.  Admissibility permits voting only and does not finalize or release anything before
the close-record decision.

\paragraph{Hash reconciliation of close versions.}
For \(T\in\{\mathsf{cut},\mathsf{frontiers}\}\), a replica retains canonical preimages indexed by
hash.  A requester that possesses base hash \(a\) and needs target hash \(b\) sends
\[
\ReconReq(T,e,a,b).
\]
A responder may answer only if it possesses both canonical preimages.  For cuts it sends
\[
\ReconDelta(\mathsf{cut},e,a,b,\Delta^+,\Delta^-),
\qquad
\Delta^+=C_b\setminus C_a,
\quad
\Delta^-=C_a\setminus C_b.
\]
The requester reconstructs
\[
C'=(C_a\setminus\Delta^-)\cup\Delta^+,
\]
checks the canonical target hash, validates \(\admissible(\CloseCut(e),b,C')\), and then caches
\(\cutVersion[e,b]=C'\).  Reconstructing a historical cut does not alter its current report store, vote store, or derived visibility.

For frontier maps the responder sends sparse per-account replacements
\[
\ReconDelta(\mathsf{frontiers},e,a,b,
\{(x,f_a(x),f_b(x)):f_a(x)\neq f_b(x)\}).
\]
The requester checks every expected old frontier, replaces only the named account entries, obtains any
missing blocks through ordinary block gossip, validates the affected chains and slot evidence, recomputes
\[
b=\hash(\mathsf{RAI/CloseRecord}\parallel e\parallel\closeHash[e-1]
\parallel\operatorname{canon}(\Phi_b)),
\]
and caches \(\frontierVersion[e,b]=\Phi_b\).  A replacement may move a pre-decision candidate frontier
forward, backward, or sideways; it never deletes a block from the immutable block store.  A responder
that lacks the base or target sends \(\ReconMiss(T,e,a,b)\), after which the requester asks another
replica or obtains a complete canonical version.  Reconciliation messages are unauthoritative transport
metadata: acceptance always requires recomputing the target hash and applicable admissibility validation.

A correct replica retains every close version for which it votes until the logical close instance
decides.  Decided cut and frontier versions, or delta paths from retained checkpoints, remain available
from correct archival replicas so that joining replicas can reconstruct the complete close history.

There is no distinguished validator proposer and no explicit proposal object in the election core.
For a slot election, the client or account owner creates and authorizes a candidate block \(B\); the
owner or any relay may gossip it, and a replica computes \(b=\hash(B)\) only after validating that
authorization and the block.  Validators do not synthesize successor blocks.

A correct replica constructs its fresh close-cut version from its current report-and-vote visibility set after
locally deriving a valid joint report quorum.  Its fresh close-cut hash is
\[
\FreshCloseHash_{P_j}(\CloseCut(e))
=
\hash(\mathsf{RAI/CloseCut}\parallel e\parallel\operatorname{canon}(C_{e,j})),
\]
and it retains the canonical set under that hash.

After the decided cut is drained, a correct replica first derives every global fast or final slot
certificate available for elections in \(\closeCut[e]\) and runs \(\FinalizeChain\) for each such
certified block in canonical account-chain order.  It then constructs a fresh close-record version from
the immutable \(\mathsf{CloseInput}_e\): a close-local replay starts from \(\closeState[e-1]\), applies
the resolver only to elections in the decided cut, and performs validated ancestry closure without
reading ordinary finalizations outside the frozen input.  This derives a complete canonical frontier
map \(\Phi_{e,j}\).  Define
\[
\FreshCloseHash_{P_j}(\CloseRecord(e))
=
\hash(\mathsf{RAI/CloseRecord}\parallel e\parallel\closeHash[e-1]
\parallel\operatorname{canon}(\Phi_{e,j})),
\]
and retain the canonical frontier map under that hash.
Fresh construction is used only when the immediately preceding close round is positively proved dead
or when starting round \(0\).  A live carried hash overrides fresh preference.

For \(r>0\), correct replica \(P_j\) chooses the first-vote value for \(I_r\) as follows:
\[
\func{NextCloseHash}_{P_j}(I,r)=
\begin{cases}
\FreshCloseHash_{P_j}(I),
  & \CloseDead_{P_j}(I_{r-1}),\\
h,
  & \LiveCarry_{P_j}(I_{r-1},h),\\
\bot,
  & \text{otherwise.}
\end{cases}
\]
For round \(0\),
\[
\func{NextCloseHash}_{P_j}(I,0)=
\FreshCloseHash_{P_j}(I).
\]
Death evidence takes precedence if the replica knows both a non-timeout notarization certificate and
a proof that the source round is dead.  If the result is \(\bot\), the replica does not start the
successor round.

A correct replica never abandons a carried hash merely because:
\begin{itemize}
\item its fresh close hash changes;
\item it learns additional valid close evidence;
\item another admissible hash becomes freshly preferred;
\item its local timer expires; or
\item it has already entered another local protocol state.
\end{itemize}
A carry may be abandoned only after a valid death proof for the round that created that carry.  Thus
\(\text{carry overrides fresh preference}\).

For a fixed account slot \((a,s)\), retries are sequential.  A replica does not start or first-vote in a later
retry \(\Slot(a,s,e')\) while some earlier election \(\Slot(a,s,e_0)\), \(e_0<e'\), remains pending.
A later same-account-slot retry may start only after the certified close record has been reconstructed and installed in the
normative order below, the resulting close state records \(\ReleasedState(\pi)\) for the earlier
election, and \(\finalizedSlot[a,s]=\bot\).  If the earlier block is finalized either by an ordinary
certificate or by the close record, no same-account-slot retry is started.  Adjacent same-account-slot retries remain
constrained by shared-committee vote locks until certified release, and disjoint future elections are
constrained by the governing certified finalized frontier.

Local completion is
\begin{center}
\begin{adjustbox}{max width=\columnwidth}
\(\displaystyle
\begin{aligned}
\Complete(\Slot(a,s,e),b)
\iff{}& \exists B:b=\hash(B)\wedge\NotarCert(\Slot(a,s,e),b)\in\mathsf{pool}\\
&\wedge\blockMap[\Slot(a,s,e),b]=B\wedge b\neq\timeoutVal_{\Slot(a,s,e)}\\
&\wedge\slotOf(B)=(a,s)\wedge\ParentComplete(B,a,s)\wedge\func{AncestorsConsistent}(B)\wedge\PayloadAvailable(B)\\
&\wedge\ValidBlock(B)\wedge\finalizedSlot[a,s]\in\{\bot,b\}.
\end{aligned}
\)
\end{adjustbox}
\end{center}
A non-genesis block enters \(\completeTree\) only by completion, and only if the same block
hash is not already present:
\begin{center}
\begin{adjustbox}{max width=\columnwidth}
\(\displaystyle
\frac{\Complete(\Slot(a,s,e),b)\quad\blockMap[\Slot(a,s,e),b]=B\quad b=\hash(B)\quad\neg\CompleteBlock(b)}
{\completeTree\gets\completeTree\cup\{B\}}.
\)
\end{adjustbox}
\end{center}
No message is broadcast solely because a block was inserted.

Epoch membership is not stored on every block.  Instead, a fresh close-record version derives one
frontier for every account and hashes the canonical frontier map together with the preceding decided
close-record hash:
\begin{center}
\begin{adjustbox}{max width=\columnwidth}
\(\displaystyle
\begin{aligned}
\Phi_{e,j}
    &=\operatorname{sort}\{(a,f_{e,j}(a)):a\in\mathsf{Accounts}\},\\
R^{\mathsf{ledger}}_{e,j}
    &=\hash(\operatorname{canon}(\Phi_{e,j})),\\
h_{F,j}
    &=\hash(\mathsf{RAI/CloseRecord}\parallel e\parallel\closeHash[e-1]
       \parallel\operatorname{canon}(\Phi_{e,j})).
\end{aligned}
\)
\end{adjustbox}
\end{center}
The decided close cut is supporting validation context for deriving \(\Phi_{e,j}\), but it is not an
additional field in the close-record value and is not itself the durable ledger-consensus object.  The
immutable \(\mathsf{CloseInput}_e\) consists of the preceding certified close, the preceding decided
frontiers, and that decided cut.  A close-record version is replica-relative before the logical close
instance decides; once \(\closeDecision[\CloseRecord(e)]\) is set, write the decided frontier map
without the replica subscript as \(\Phi_e\).  Blocks, reports, and votes are obtained through ordinary
gossip, while the frontier map itself is reconstructed by hash reconciliation when necessary.

Because every block contains its predecessor hash, each frontier transitively commits to its entire
account chain.  The validator recomputes every block hash and owner signature, checks account and
sequence continuity, replays every debit and credit, verifies finalized send-before-receive dependencies,
checks destination ownership and unique send consumption, applies every representative change, and
rejects any chain with a missing or invalid predecessor.  The replay begins from the preceding certified
close state and applies only blocks admitted through the frozen \(\mathsf{CloseInput}_e\); ordinary
finalizations from elections outside the decided epoch-\(e\) cut do not affect this replay.  Every receive
must resolve to a still-pending finalized send addressed to the receiving account.  It also requires
every block on the preceding certified account chains to remain, every balance to remain non-negative,
and the supply invariant to hold.
The resulting \(\bal_e\), \(\repOf_e\), \(\pendingSend_e\), and \(\receivedSend_e\) maps are part of
the reconstructed \(\closeState[e]\).  The canonical hash of the complete frontier map is
\(R^{\mathsf{ledger}}_e\), while the close-record hash additionally commits to epoch \(e\) and
\(\closeHash[e-1]\).

Let \(\frontier[e,a]\) be the decided frontier of account \(a\) through epoch \(e\), with
\(\frontier[e,a]=\frontier[e-1,a]\) for accounts that do not move in epoch \(e\).  The blocks of
account \(a\) belonging to epoch \(e\) are the ordered account-chain segment after
\(\frontier[e-1,a]\) and up to \(\frontier[e,a]\):
\begin{center}
\begin{adjustbox}{max width=\columnwidth}
\(\displaystyle
\EpochBlocks(e,a)=
(\frontier[e-1,a],\frontier[e,a]] .
\)
\end{adjustbox}
\end{center}
Only the frontier block is named in \(\Phi_e\); every intermediate block is committed through the
predecessor-hash chain and its epoch membership is obtained from account-tree order.

For \(q=\Slot(a,s,e)\), the decided ledger includes candidate \(b\) exactly when
\(\func{LedgerIncludes}_e(q,b)\).  Every such block and every ancestor on its certified account chain
is finalized.  Every other candidate of \(q\) is released.  If no candidate of \(q\) occurs on the
certified ledger, then \(q\) is released.  Every epoch-\(e\) slot election outside the certified cut
is released by close exclusion.  Before decision, ordinary-finalized and selected are evidence
classifications only; after certification the durable distinction is whether a block occurs on a chain
committed by the decided frontier map.

Once the logical close-record instance decides \(h_F\), a replica installs the corresponding
reconstructed frontier version atomically in the following normative order: (1) locally derive the
deciding fast or final certificate from signed close-record votes; (2) reconstruct \(\Phi_e\) for \(h_F\) by a
retained version, reconciliation delta, or complete version transfer; (3) verify
\[
h_F=\hash(\mathsf{RAI/CloseRecord}\parallel e\parallel\closeHash[e-1]
\parallel\operatorname{canon}(\Phi_e));
\]
(4) reconstruct and historically validate the decided close-cut version and locally derive every required
slot certificate, timeout, or conflict from signed votes; (5) obtain and validate every block on every account
chain from genesis through the claimed frontiers; (6) verify the frontier classifications, ancestry
closure, ledger transitions, and \(R^{\mathsf{ledger}}_e=\hash(\operatorname{canon}(\Phi_e))\); (7)
run \(\FinalizeChain\) in account-chain order through every certified frontier; and only then (8)
derive and install \(\closeState[e]\), \(\frontier[e,\cdot]\),
\(\ledgerRoot[e]\gets R^{\mathsf{ledger}}_e\), and \(\closeHash[e]\gets h_F\).  No retry or
future-vote decision may read the new release state before all eight steps complete.

A joining replica may wait for the next live close to learn the latest synchronization target, but that
observation is not a trust assumption.  It starts from genesis, obtains the durable blocks, reports, and
votes, and replays these steps for every epoch in order.  It reconstructs each decided cut and frontier
preimage, validates every block and state transition, locally derives each deciding vote quorum,
recomputes every close-record hash from \(\closeHash[e-1]\) and \(\Phi_e\), and accepts epoch \(e\)
only when the recomputed hash is the value decided by the close-record votes.  The latest validated
close-record hash is the current consensus checkpoint: its predecessor-hash chain and account
frontiers transitively commit to every block on the certified ledger chains.  Thus the replica derives
the live committee internally and bootstraps the entire ledger without trusting an archive's derived
state or an externally supplied committee snapshot.

If \(\FirstCert(\Slot(a,s,e),b)\in\mathsf{pool}\),
\(\FinalCert(\Slot(a,s,e),b)\in\mathsf{pool}\), or a decided ledger-state root includes \(b\) on a
certified account chain, with \(b\neq\timeoutVal_{\Slot(a,s,e)}\) and \(\CompleteBlock(b)\), then
\(\FinalizeChain(b)\) runs.  It first requires
\(\func{AncestorsConsistent}(B)\wedge\StablePrefix(B)\), then walks the complete parent chain in
ancestor-to-descendant order, \(b_0=\mathsf{genesis}_{\accountOf(B_k)},b_1,\ldots,b_k=b\).  For each
\(i>0\), with \(b_i=\hash(B_i)\) and \(\slotOf(B_i)=(a_i,s_i)\), it first asserts
\(\finalizedSlot[a_i,s_i]\in\{\bot,b_i\}\).  A violation is a safety fault: a finality proof has appeared
for a chain conflicting with an already finalized account slot.  If \(\finalizedSlot[a_i,s_i]=\bot\),
it validates \(B_i\) against the current finalized ledger state, applies its balance, send, receive, and
representative transition, and then sets \(\finalizedSlot[a_i,s_i]\gets b_i\).  If
\(\finalizedSlot[a_i,s_i]=b_i\), it does nothing.  The derived finalized ledger state is therefore
updated before a later block is validated.  Timeout and release outcomes never invoke
\(\FinalizeChain\).

Normative invariants: enabled voting; first-vote uniqueness; certified release for conflicting slot
vote history; timeout locality; separation of notarization from slot finality; no gaps in
\(\completeTree\); at most one complete-tree entry per block hash; certified per-account epoch
frontiers; a canonical frontier map and hash-chained close-record value over every account frontier; complete validation of every
account chain from genesis through its frontier; valid owner signatures; non-negative balances; finalized
send-before-receive dependencies; destination-only and at-most-once send consumption; deterministic
representative updates; supply conservation; monotonic \(\finalizedSlot\); ancestor finality by
ordinary fast or final certificates or by certified ledger inclusion across a certified-stable ancestor
prefix; release of every epoch election omitted from the certified ledger; close-cut boundedness after a
certified close-cut hash; a locally derivable joint report quorum and valid reconstructed cut for every non-timeout close-cut vote;
mandatory validation under \(\closeState[e-2]\) for every epoch-\(e\) vote; deterministic stake-weighted committee
derivation from certified balances, representatives, close state, and close-record hash; uniqueness and irrevocability of the certified close
record; and at most one open epoch and at most one closing epoch, with older epochs closed.

\section{Visibility and Close Cuts}

\begin{center}
\begin{adjustbox}{max width=\columnwidth}
\(\displaystyle
\begin{aligned}
\Reportable_j(\Slot(a,s,e))\iff{}& \started[\Slot(a,s,e)]\wedge
\bigl(\neg\finished(\Slot(a,s,e))\vee
\exists Q\in\Committees(\Slot(a,s,e)),\ \exists b\neq\timeoutVal_{\Slot(a,s,e)}:\\
&\qquad b\in\func{Support}(P_j,Q,\Slot(a,s,e))\wedge
\neg\released_{P_j}(\Slot(a,s,e),b)\bigr).
\end{aligned}
\)
\end{adjustbox}
\end{center}
A correct report therefore retains every unreleased non-timeout slot-vote obligation of its signer,
even if that signer has already observed a replica-local terminal outcome.

For every signer and epoch, \(\reportStore[P_i,e]\) retains every distinct valid signed report learned
for that signer.  Report equivocation and eligibility are derived objectively as
\[
\begin{aligned}
\ReportEquivocated(P_i,e)\iff{}&\exists R\neq R':
\Report(e,R,\sigma),\Report(e,R',\sigma')\in\reportStore[P_i,e],\\
\EligibleReport(P_i,e,R)\iff{}&\Report(e,R,\sigma)\in\reportStore[P_i,e]
\wedge\neg\ReportEquivocated(P_i,e).
\end{aligned}
\]
The derived value \(report[P_i,e]\) is the unique body \(R\) satisfying
\(\EligibleReport(P_i,e,R)\), and is \(\bot\) otherwise.  Thus all signed reports are retained,
but an equivocating signer contributes no weight to report quorums or report-derived visibility.

A joint report proof \(\rho\) contains, for every close committee, distinct eligible signed reports whose
combined signer weight crosses the close barrier:
\begin{center}
\begin{adjustbox}{max width=\columnwidth}
\(\displaystyle
\JointReportProof(e,\rho)\iff
\forall Q\in\Committees(\CloseCut(e,0)):\
\SignerWeight_Q(\{P_i\in Q:\exists R_i,\ \Report(e,R_i,\sigma_i)\in\rho_Q
\wedge\EligibleReport(P_i,e,R_i)\})\ge W_Q-F_Q.
\)
\end{adjustbox}
\end{center}
All reports in \(\rho\) must be individually valid and eligible.  In particular, no signer for which
report equivocation is known contributes weight to the proof.  The account slots forced by such a proof are
\begin{center}
\begin{adjustbox}{max width=\columnwidth}
\(\displaystyle
\ReportVisible(e,\rho)=
\{(a,s):\forall Q\in\Committees(\Slot(a,s,e)),\
\SignerWeight_Q(\{P_i\in Q:\exists R_i,\ \Report(e,R_i,\sigma_i)\in\rho_Q
\wedge\EligibleReport(P_i,e,R_i)\wedge\Slot(a,s,e)\in R_i\})>F_Q\}.
\)
\end{adjustbox}
\end{center}

\begin{breakablealgorithm}
\caption{\(\UpdateVisible(e)\)}
\begin{algorithmic}[1]
\State $V_{\mathrm{reports}}\gets\{(a,s):\forall Q\in\Committees(\Slot(a,s,e)),\ \SignerWeight_Q(\{P_i\in Q:\Slot(a,s,e)\in report[P_i,e]\})>F_Q\}$
\State $V_{\mathrm{votes}}\gets\{(a,s):\exists Q\in\Committees(\Slot(a,s,e)):\mathsf{pool}\text{ contains valid votes of distinct-signer weight }>F_Q\text{ in }Q\text{ for }\Slot(a,s,e)\}$
\State $V\gets V_{\mathrm{reports}}\cup V_{\mathrm{votes}}$
\If{$V\neq visible[e]$}
    \State $visible[e]\gets V$
    \If{$state[e]=\mathsf{closing}$ and a cached valid joint report proof $\rho$ exists}
        \State $C\gets\{\Slot(a,s,e):(a,s)\in visible[e]\cup\ReportVisible(e,\rho)\}$
        \State $h_C\gets\hash(\mathsf{RAI/CloseCut}\parallel e\parallel\operatorname{canon}(C))$
        \State $\cutVersion[e,h_C]\gets C$
    \EndIf
\EndIf
\end{algorithmic}
\end{breakablealgorithm}

Reports contain pending slot elections and every unreleased non-timeout slot-vote obligation of
the reporter.  Report and vote visibility are evaluated per weighted committee, using the committees of
the slot election whose visibility is being tested.  Report-derived visibility is recomputed from the
currently eligible reports and may shrink when report equivocation is learned; vote-derived visibility
remains monotone because valid votes are persistent.  A replica's current close cut is
\(C=\{\Slot(a,s,e):(a,s)\in visible[e]\}\), and its close-cut hash is
\[
\hash(\mathsf{RAI/CloseCut}\parallel e\parallel\operatorname{canon}(C)),
\]
not a hash of reports and not the durable epoch close-record hash.  A slot election enters current
visibility either through report signer weight greater than \(F_Q\) in every committee
\(Q\in\Committees(\Slot(a,s,e))\), or through visible valid-vote signer weight greater than \(F_Q\)
in some such committee.  Before a non-timeout close-cut vote, the replica locally selects a valid joint
report proof \(\rho\), requires the canonical cut to contain both its current visible set and every
account slot in \(\ReportVisible(e,\rho)\), validates a local \(\func{StartWitness}_e(q)\) for every
\(q\in C\), and stores \(\cutVersion[e,h_C]=C\).  No cut, proof, witness map, or certificate is
broadcast; only the close vote containing \(h_C\) is disseminated.

If \(\closeDecision[\CloseCut(e)]=(r,h_C)\), the replica reconstructs and historically validates
\[
C=\cutVersion[e,h_C]
\]
by retained state or hash reconciliation and sets \(\closeCut[e]=C\).  This fixes
\(\mathsf{CloseInput}_e=(\closeHash[e-1],\frontier[e-1,\cdot],C)\).  Its current visibility may differ
from \(C\); reconstructing the decided historical version does not remove reports or votes and does not
overwrite the visibility derived from those stores.  For every \(\Slot(a,s,e)\in C\), each replica
derives a fresh close-record frontier map by close-local replay from the preceding certified state,
using the valid slot evidence currently known only for elections in \(C\) and the validated block tree.
Different correct replicas may initially derive different admissible hashes because their evidence for
the frozen elections differs.  Newly gossiped atoms may change fresh preference but cannot invalidate a
previously validated historical version or override a live carry.  Replicas vote on
\[
h_F=\hash(\mathsf{RAI/CloseRecord}\parallel e\parallel\closeHash[e-1]
\parallel\operatorname{canon}(\Phi_e)).
\]
Only the value stored in \(\closeDecision[\CloseRecord(e)]\) becomes \(\closeHash[e]\); the close-cut
hash and local slot-drain outcomes do not.

\section{Main Algorithms}

\begin{breakablealgorithm}
\caption{RAI initialization for replica $P_j$}
\begin{algorithmic}[1]
\State $e_{\mathrm{open}}\gets0$; $state[-1]\gets\mathsf{closed}$; $state[e_{\mathrm{open}}]\gets\mathsf{open}$; let $G$ be the genesis epoch state
\State $\closeState[-1]\gets G$; $\closeHash[-1]\gets\hash(G)$; $\committee_{\mathsf{gen}}\gets\deriveCommittee(\mathsf{genesis},G)$; $\committee[\cdot]\gets\bot$
\State $\completeTree\gets\{\mathsf{genesis}_a:a\in\mathsf{Accounts}\}$; $\frontier[-1,a]\gets\mathsf{genesis}_a$ for every account $a$; derive $\bal_{-1}$, $\repOf_{-1}$, $\pendingSend_{-1}$, and $\receivedSend_{-1}$ from $G$
\State $\ledgerRoot[-1]\gets\hash(\operatorname{sort}\{(a,\frontier[-1,a]):a\in\mathsf{Accounts}\})$
\State $\finalizedSlot[\cdot,\cdot]\gets\bot$
\State $epochStart[\cdot]\gets\bot$; $epochStart[e_{\mathrm{open}}]\gets\func{clock}()$; $T_{\mathrm{election}}[\cdot]\gets\bot$
\State $outcome[\cdot]\gets\pending$; $\closeDecision[\cdot]\gets\bot$; $\closeRound[\cdot]\gets0$; $notarized[\cdot,\cdot,\cdot]\gets\{\}$; $secondLook[\cdot,\cdot,\cdot]\gets\{\}$
\State $firstVote[\cdot,\cdot,\cdot]\gets\bot$; $finalVoted[\cdot,\cdot,\cdot]\gets\false$; $\votedFor[\cdot,\cdot,\cdot]\gets\false$; $\blockMap[\cdot,\cdot]\gets\bot$
\State $visible[e_{\mathrm{open}}]\gets\{\}$; $\reportStore[\cdot,e_{\mathrm{open}}]\gets\{\}$; $\closeHash[e_{\mathrm{open}}]\gets\bot$; $\ledgerRoot[e_{\mathrm{open}}]\gets\bot$; $\closeState[e_{\mathrm{open}}]\gets\bot$; $\cutVersion[e_{\mathrm{open}}]\gets\{\}$; $\frontierVersion[e_{\mathrm{open}}]\gets\{\}$; $\closeCut[e_{\mathrm{open}}]\gets\bot$
\end{algorithmic}
\end{breakablealgorithm}

\begin{breakablealgorithm}
\caption{\(\StartClosing(k)\)}
\begin{algorithmic}[1]
\State \textbf{require} $state[k]=\mathsf{open}$ and $state[k-1]=\mathsf{closed}$ and $\closeState[k-1]\neq\bot$
\State $state[k]\gets\mathsf{closing}$
\Statex \Comment{The epoch-\(k\) frontier input will be fixed objectively by the decided cut and \(\frontier[k-1,\cdot]\), before any epoch-\(k+1\) election may enter epoch \(k\)'s close record.}
\State $\closeDecision[\CloseCut(k)]\gets\bot$; $\closeRound[\CloseCut(k)]\gets0$; $\closeDecision[\CloseRecord(k)]\gets\bot$; $\closeRound[\CloseRecord(k)]\gets0$
\State $R_j[k]\gets\{\Slot(a,s',k):a\in\mathsf{Accounts},\ s'\ge1,\ \Reportable_j(\Slot(a,s',k))\}$
\State $\reportStore[P_j,k]\gets\reportStore[P_j,k]\cup\{\Report(k,R_j[k],\sigma_j)\}$; broadcast $\Report(k,R_j[k],\sigma_j)$
\State $e_{\mathrm{open}}\gets k+1$; $state[e_{\mathrm{open}}]\gets\mathsf{open}$; $epochStart[e_{\mathrm{open}}]\gets\func{clock}()$
\State $visible[e_{\mathrm{open}}]\gets\{\}$; $\reportStore[\cdot,e_{\mathrm{open}}]\gets\{\}$; $\closeHash[e_{\mathrm{open}}]\gets\bot$; $\ledgerRoot[e_{\mathrm{open}}]\gets\bot$; $\closeState[e_{\mathrm{open}}]\gets\bot$; $\cutVersion[e_{\mathrm{open}}]\gets\{\}$; $\frontierVersion[e_{\mathrm{open}}]\gets\{\}$; $\closeCut[e_{\mathrm{open}}]\gets\bot$
\State load the certified finalized frontiers and release evidence from $\closeState[k-1]$ as durable validation state
\end{algorithmic}
\end{breakablealgorithm}

\begin{breakablealgorithm}
\caption{\(\FinishClosing(k)\)}
\begin{algorithmic}[1]
\State \textbf{require} $state[k]=\mathsf{closing}$ and $state[k-1]=\mathsf{closed}$
\State $state[k]\gets\mathsf{closed}$; $\func{clearEpochVars}(k)$
\end{algorithmic}
\end{breakablealgorithm}

\begin{breakablealgorithm}
\caption{\(\func{ObtainCloseVersion}_{P_j}(I,h)\)}
\begin{algorithmic}[1]
\If{$I=\CloseCut(e)$ and $\cutVersion[e,h]=C$ and $\admissible(I,h,C)$}
    \State \Return $C$
\ElsIf{$I=\CloseRecord(e)$ and $\frontierVersion[e,h]=\Phi$ and $\admissible(I,h,\Phi)$}
    \State \Return $\Phi$
\EndIf
\ForAll{locally known base hashes $a$ for the same close instance}
    \State send $\ReconReq(T,e,a,h)$, where $T=\mathsf{cut}$ or $\mathsf{frontiers}$
    \If{a valid $\ReconDelta(T,e,a,h,\delta)$ is received}
        \State apply $\delta$ to the retained preimage of $a$
        \State recompute exactly $h$ and run the applicable admissibility checks
        \If{validation succeeds}
            \State cache and \Return the reconstructed version
        \EndIf
    \EndIf
\EndFor
\State retrieve a complete canonical version for $h$
\State recompute exactly $h$ and run the applicable admissibility checks
\State cache and \Return the version
\end{algorithmic}
\end{breakablealgorithm}

\begin{breakablealgorithm}
\caption{Start close round \(I_r\) at correct replica \(P_j\)}
\begin{algorithmic}[1]
\State \textbf{require} $\closeDecision[I]=\bot$
\If{$r=0$}
    \State $x\gets\FreshCloseHash_{P_j}(I)$
\ElsIf{$\CloseDead_{P_j}(I_{r-1})$}
    \State $x\gets\FreshCloseHash_{P_j}(I)$
\ElsIf{$\exists h:\LiveCarry_{P_j}(I_{r-1},h)$}
    \State let $h$ be the locally live carried hash
    \State $S\gets\func{ObtainCloseVersion}_{P_j}(I,h)$
    \State $x\gets h$
\Else
    \State \Return without starting round $r$
\EndIf
\State initialize the committee-local election-kernel state for $I_r$
\State $\closeRound[I]\gets\max(\closeRound[I],r)$
\State first-vote $x$ in every applicable close committee
\end{algorithmic}
\end{breakablealgorithm}

When \(x\) is a carried hash, the replica waits to reconstruct and validate its canonical version rather than
replacing \(x\) with a timeout first vote.  The replica does not issue a timeout first vote in \(I_r\)
merely because the carried version is still being reconstructed; equivalently,
\(\text{no timeout first vote while the required carried version is being reconstructed}\).

After the carried first vote has been issued, the committee-local second-look, notarization,
timeout-notarization, fast, and final rules remain enabled.

Here \(\func{clearEpochVars}(e)\) deletes only volatile caches.  It never deletes durable safety
evidence, including \(\closeHash[e]\), \(\ledgerRoot[e]\), \(\closeState[e]\), \(\closeCut[e]\),
\(\committee[e]\), decided \(\cutVersion\) and \(\frontierVersion\) entries, timeout, conflict,
close-exclusion, or certified-ledger evidence, local unreleased votes, \(\completeTree\),
\(\frontier[\cdot,\cdot]\), or \(\finalizedSlot\).

\begin{breakablealgorithm}
\caption{RAI epoch loop for replica $P_j$}
\begin{algorithmic}[1]
\State \textbf{while} true \textbf{wait until either}
\Statex
\State $state[e_{\mathrm{open}}]=\mathsf{open}$ and $state[e_{\mathrm{open}}-1]=\mathsf{closed}$ and $\func{clock}()>epochStart[e_{\mathrm{open}}]+\Delta_{\mathrm{timeout}}$ $\Rightarrow$
\State \quad $k\gets e_{\mathrm{open}}$; $\StartClosing(k)$
\Statex
\State received valid $\Report(k,R_i[k],\sigma_i)$ from $P_i$ with $\Report(k,R_i[k],\sigma_i)\notin\reportStore[P_i,k]$ $\Rightarrow$
\State \quad $\reportStore[P_i,k]\gets\reportStore[P_i,k]\cup\{\Report(k,R_i[k],\sigma_i)\}$; recompute the derived $report[P_i,k]$; $\UpdateVisible(k)$
\Statex
\State $state[k]=\mathsf{closing}$ and $\forall Q\in\Committees(\CloseCut(k,0)):
\SignerWeight_Q(\{P_i\in Q:report[P_i,k]\neq\bot\})\ge W_Q-F_Q$ and $\closeDecision[\CloseCut(k)]=\bot$ and $\closeRound[\CloseCut(k)]=0$ $\Rightarrow$
\State \quad assemble a valid joint report proof $\rho_k$ from the cached reports; $\UpdateVisible(k)$
\State \quad $C\gets\{\Slot(a,s,k):(a,s)\in visible[k]\cup\ReportVisible(k,\rho_k)\}$
\State \quad $h_C\gets\hash(\mathsf{RAI/CloseCut}\parallel k\parallel\operatorname{canon}(C))$; $\cutVersion[k,h_C]\gets C$
\State \quad start close round $\CloseCut(k,0)$ using $\func{NextCloseHash}_{P_j}(\CloseCut(k),0)$
\Statex
\State $\closeDecision[\CloseCut(k)]=(r,h_C)$ and $\closeCut[k]=\bot$ $\Rightarrow$
\State \quad locally derive the deciding $\FirstCert(\CloseCut(k,r),h_C)$ or $\FinalCert(\CloseCut(k,r),h_C)$ from signed votes
\State \quad $C\gets\func{ObtainCloseVersion}_{P_j}(\CloseCut(k),h_C)$; validate $\admissibleHistorical(\CloseCut(k),h_C,C)$
\State \quad $\closeCut[k]\gets C$; fix $\mathsf{CloseInput}_k=(\closeHash[k-1],\frontier[k-1,\cdot],C)$
\State \quad for every $q\in C$: derive its historical start obligation from the certified cut and, if $T_{\mathrm{election}}[q]=\bot$, set $T_{\mathrm{election}}[q]\gets\func{clock}()$
\State \quad \textbf{wait until} $\forall q=\Slot(a,s,k)\in C:\finished(q)$
\State \quad derive every currently available global fast or final certificate for slots in $C$ and run $\FinalizeChain$ for each certified block in canonical account-chain order
\State \quad derive $\Phi_k$ by close-local replay from $\closeState[k-1]$, using current valid evidence only for $q\in C$; ignore ordinary finalizations outside the frozen $\mathsf{CloseInput}_k$
\State \quad $R^{\mathsf{ledger}}_k\gets\hash(\operatorname{canon}(\Phi_k))$
\State \quad $h_F\gets\hash(\mathsf{RAI/CloseRecord}\parallel k\parallel\closeHash[k-1]\parallel\operatorname{canon}(\Phi_k))$
\State \quad $\frontierVersion[k,h_F]\gets\Phi_k$
\State \quad start close round $\CloseRecord(k,0)$ using $\func{NextCloseHash}_{P_j}(\CloseRecord(k),0)$
\Statex
\State $\closeDecision[\CloseRecord(k)]=(r,h_F)$ and $\closeHash[k]=\bot$ $\Rightarrow$
\State \quad locally derive the deciding $\FirstCert(\CloseRecord(k,r),h_F)$ or $\FinalCert(\CloseRecord(k,r),h_F)$ from signed votes
\State \quad $\Phi_k\gets\func{ObtainCloseVersion}_{P_j}(\CloseRecord(k),h_F)$; validate $\admissible(\CloseRecord(k),h_F,\Phi_k)$
\State \quad $R^{\mathsf{ledger}}_k\gets\hash(\operatorname{canon}(\Phi_k))$; require $h_F=\hash(\mathsf{RAI/CloseRecord}\parallel k\parallel\closeHash[k-1]\parallel\operatorname{canon}(\Phi_k))$
\State \quad obtain and validate every account chain from genesis through the frontiers in $\Phi_k$
\State \quad replay and verify every owner signature, balance, send, receive, and representative transition; finalize every block on those chains and derive released elections by omission
\State \quad install $\closeState[k]$ with $\bal_k$, $\repOf_k$, $\pendingSend_k$, and $\receivedSend_k$, plus $\frontier[k,\cdot]$ and $\ledgerRoot[k]\gets R^{\mathsf{ledger}}_k$; $\closeHash[k]\gets h_F$; $\committee[k]\gets\deriveCommittee(k,\closeHash[k],\closeState[k])$
\State \quad $\FinishClosing(k)$
\end{algorithmic}
\end{breakablealgorithm}

\begin{breakablealgorithm}
\caption{Generic slot-election loop for $\id=\Slot(a,s,e)$ at replica $P_j$}
\begin{algorithmic}[1]
\State \textbf{while} $\neg\finished(\Slot(a,s,e))$ and $state[e]\neq\mathsf{closed}$ \textbf{wait until either}
\Statex
\State $\id=\Slot(a,s,e)$ and $\exists b,B:b=\hash(B)\wedge\Complete(\Slot(a,s,e),b)\wedge\blockMap[\Slot(a,s,e),b]=B$ $\Rightarrow$
\If{$\neg\CompleteBlock(b)$}
    \State $\completeTree\gets\completeTree\cup\{B\}$
\EndIf
\State \quad $B_{\mathrm p}[\id]\gets b$
\State \quad $\mathcal Q_f\gets\{Q\in\Committees(\id):P_j\in Q\wedge\neg\finalVoted(P_j,\id,Q)\wedge \func{Support}(P_j,Q,\id)\subseteq\{b\}\}$
\If{$\enabled(\id)$ and $\mathcal Q_f=\{Q\in\Committees(\id):P_j\in Q\}$ and $\safeToVote_{P_j}(\Slot(a,s,e),b)$}
    \State $\votedFor[P_j,\Slot(a,s,e),b]\gets\true$
    \State for every $Q\in\mathcal Q_f$: $finalVoted[P_j,\id,Q]\gets\true$; logically emit $\FinalVote(\id,Q,b,\sigma_j)$
    \State broadcast one unscoped final vote when it represents all of those scoped votes, otherwise broadcast the scoped final votes
\EndIf
\Statex
\State the pool contains a certificate for $(t,\id,x)$ with $\Terminal(t,\id,x)$ $\Rightarrow$
\If{$outcome[\id]=\pending$}
    \State $outcome[\id]\gets\certOutcome(t,x)$
\EndIf
\If{$\id=\Slot(a,s,e)$ and $(\FirstCert(\id,x)\in\mathsf{pool}\vee\FinalCert(\id,x)\in\mathsf{pool})$ and $x\neq\timeoutVal_\id$ and $\CompleteBlock(x)$}
    \State $\FinalizeChain(x)$
\EndIf
\Statex
\State $\exists Q\in\Committees(\id):P_j\in Q\wedge\neg\firstVoted(P_j,\id,Q)$ and locally derived, reconstructed, or learned the required data for value $x$ for $\id$ $\Rightarrow$
\State \quad choose such a $Q$; \textbf{require} $\enabled(\id)$ and $\neg\finalVoted(P_j,\id,Q)$
\If{$\id=\Slot(a,s,e)$ and $x=b=\hash(B)$ for the received block $B$}
    \State $\blockMap[\Slot(a,s,e),b]\gets B$
\EndIf
\State \quad \textbf{require} $\admissible(\id,x)$
\If{$\id=\Slot(a,s,e)$}
    \State \quad \textbf{require} $\safeToVote_{P_j}(\Slot(a,s,e),x)$
\EndIf
\If{$T_{\mathrm{election}}[\id]=\bot$}
    \State $T_{\mathrm{election}}[\id]\gets\func{clock}()$
\EndIf
\State \quad $firstVote[P_j,\id,Q]\gets x$
\If{$\id=\Slot(a,s,e)$}
    \State \quad $\votedFor[P_j,\Slot(a,s,e),x]\gets\true$
\EndIf
\State \quad logically emit $\FirstVote(\id,Q,x,\sigma_j)$; an identical vote for every applicable committee may be sent once as an unscoped wire vote
\Statex
\State $\exists Q\in\Committees(\id):P_j\in Q\wedge\neg\firstVoted(P_j,\id,Q)$ and $T_{\mathrm{election}}[\id]\neq\bot$ and $\func{clock}()>T_{\mathrm{election}}[\id]+\Delta_{\mathrm{timeout}}$ $\Rightarrow$
\State \quad choose such a $Q$; \textbf{require} $\enabled(\id)$ and $\neg\finalVoted(P_j,\id,Q)$
\State \quad $firstVote[P_j,\id,Q]\gets\timeoutVal_\id$
\State \quad logically emit $\FirstVote(\id,Q,\timeoutVal_\id,\sigma_j)$; identical scoped votes may be compressed into one unscoped wire vote
\Statex
\State received a valid scoped $\FirstVote(\id,Q,x,\_)$ from $P_i$ and $firstVote[P_i,\id,Q]=\bot$ $\Rightarrow$
\If{$T_{\mathrm{election}}[\id]=\bot$}
    \State $T_{\mathrm{election}}[\id]\gets\func{clock}()$
\EndIf
\State \quad $firstVote[P_i,\id,Q]\gets x$; add the scoped vote to $\mathsf{pool}$
\Statex
\State received a valid scoped $\NotarVote(\id,Q,x,\_)$ from $P_i$ $\Rightarrow$
\State \quad add the scoped vote to $\mathsf{pool}$
\Statex
\State received a valid scoped $\FinalVote(\id,Q,x,\_)$ from $P_i$ $\Rightarrow$
\State \quad add the scoped vote to $\mathsf{pool}$
\Statex
\State $\exists Q\in\Committees(\id),\ x\in\manyVotes(Q,\id)\setminus secondLook[P_j,Q,\id]$ and $\admissible(\id,x)$ $\Rightarrow$
\State \quad choose such a $(Q,x)$; \textbf{require} $P_j\in Q$, $\enabled(\id)$, $\firstVoted(P_j,\id,Q)$, and $\neg\finalVoted(P_j,\id,Q)$
\If{$\id=\Slot(a,s,e)$}
    \State \quad \textbf{require} $\safeToVote_{P_j}(\Slot(a,s,e),x)$
\EndIf
\State \quad $secondLook[P_j,Q,\id]\gets secondLook[P_j,Q,\id]\cup\{x\}$
\If{$x\notin notarized[P_j,Q,\id]$}
    \If{$\id=\Slot(a,s,e)$}
        \State $\votedFor[P_j,\Slot(a,s,e),x]\gets\true$
    \EndIf
    \State broadcast $\NotarVote(\id,Q,x,\sigma_j)$; $notarized[P_j,Q,\id]\gets notarized[P_j,Q,\id]\cup\{x\}$
\EndIf
\Statex
\State $\exists Q\in\Committees(\id):P_j\in Q\wedge\firstVoted(P_j,\id,Q)\wedge\timeoutAllowed(Q,\id)\wedge\timeoutVal_\id\notin notarized[P_j,Q,\id]$ $\Rightarrow$
\State \quad choose such a $Q$; \textbf{require} $\enabled(\id)$ and $\neg\finalVoted(P_j,\id,Q)$
\State \quad broadcast $\NotarVote(\id,Q,\timeoutVal_\id,\sigma_j)$; $notarized[P_j,Q,\id]\gets notarized[P_j,Q,\id]\cup\{\timeoutVal_\id\}$
\Statex
\State $\id=\Slot(a,s,e)$ and $(\FirstCert(\id,b)\in\mathsf{pool}\vee\FinalCert(\id,b)\in\mathsf{pool})$ and $b\neq\timeoutVal_\id$ and $\CompleteBlock(b)$ $\Rightarrow$
\State \quad $\FinalizeChain(b)$
\Statex
\State $\id=\Slot(a,s,e)$ and $\exists Q\in\Committees(\Slot(a,s,e)):$ the pool contains valid votes of distinct-signer weight $>F_Q$ in $Q$ for $\Slot(a,s,e)$ $\Rightarrow$
\State \quad $\UpdateVisible(e)$
\end{algorithmic}
\end{breakablealgorithm}

\begin{breakablealgorithm}
\caption{Close-round processing for \(I_r\) at correct replica \(P_j\)}
\begin{algorithmic}[1]
\State use the committee-local receive, vote-counting, fast, final, and timeout-notarization rules for $I_r$
\State $\exists x\in\manyVotes(Q,I_r)\setminus secondLook[P_j,Q,I_r]$ $\Rightarrow$
\State \quad reconstruct and validate the canonical cut or frontier version for $x$
\State \quad add $x$ to $secondLook[P_j,Q,I_r]$
\If{$x\notin notarized[P_j,Q,I_r]$}
    \State broadcast $\NotarVote(I_r,Q,x,\sigma_j)$
    \State $notarized[P_j,Q,I_r]\gets notarized[P_j,Q,I_r]\cup\{x\}$
\EndIf
\Statex
\State $\exists Q\in\Committees(I_r),\ h\neq\timeoutVal_{I_r}:P_j\in Q\wedge\firstVoted(P_j,I_r,Q)\wedge\neg\finalVoted(P_j,I_r,Q)\wedge\func{Support}(P_j,Q,I_r)\subseteq\{h\}$ and the canonical version for $h$ has been reconstructed and validated $\Rightarrow$
\State \quad choose such a $(Q,h)$; \textbf{require} $\enabled(I_r)$
\State \quad $finalVoted[P_j,I_r,Q]\gets\true$; logically emit $\FinalVote(I_r,Q,h,\sigma_j)$
\State \quad when the same final-vote action is enabled in every applicable committee, broadcast one unscoped final vote; otherwise broadcast the scoped vote
\Statex
\State the local vote store derives $\FirstCert(I_r,h)$ for non-timeout $h$ $\Rightarrow$
\State \quad reconstruct and validate the canonical version for $h$
\If{$\closeDecision[I]=\bot$}
    \State $\closeDecision[I]\gets(r,h)$
\ElsIf{$\closeDecision[I].hash\neq h$}
    \State signal a safety violation
\EndIf
\State stop emitting further votes for close instance $I$
\Statex
\State the local vote store derives $\FinalCert(I_r,h)$ for non-timeout $h$ $\Rightarrow$
\State \quad reconstruct and validate the canonical version for $h$
\If{$outcome[I_r]=\pending$}
    \State $outcome[I_r]\gets\final(h)$
\EndIf
\If{$\closeDecision[I]=\bot$}
    \State $\closeDecision[I]\gets(r,h)$
\ElsIf{$\closeDecision[I].hash\neq h$}
    \State signal a safety violation
\EndIf
\State stop emitting further votes for close instance $I$
\Statex
\State the local vote store derives $\NotarCert(I_r,h)$ for non-timeout $h$ $\Rightarrow$
\If{$outcome[I_r]=\pending$}
    \State $outcome[I_r]\gets\converged(h)$
\EndIf
\If{$\closeDecision[I]=\bot$}
    \State enable round $r+1$ using $\func{NextCloseHash}_{P_j}(I,r+1)$
\EndIf
\Statex
\State $\CloseDead_{P_j}(I_r)$ $\Rightarrow$
\If{$outcome[I_r]=\pending$}
    \State $outcome[I_r]\gets\timeoutOutcome$
\EndIf
\If{$\closeDecision[I]=\bot$}
    \State enable round $r+1$ using $\func{NextCloseHash}_{P_j}(I,r+1)$
\EndIf
\end{algorithmic}
\end{breakablealgorithm}

\paragraph{Reconciliation metadata in close-vote gossip.}
A correct replica may advertise retained close-version hashes and available delta bases alongside ordinary
vote gossip.  Such advertisements are unsigned transport metadata unless authenticated by the transport
layer; they are not votes, certificates, certified values, or decisions.  A receiver accepts a
reconstructed version only after applying a delta to a retained canonical base, recomputing the exact
target hash, and running applicable admissibility validation.  Invalid, inapplicable, or unavailable deltas
are ignored, and the receiver asks another replica or obtains a complete canonical version.

\section{Safety}

Assume every weighted committee \(Q\) has Byzantine signer weight at most \(F_Q\), quorum slack
\(P_Q\ge F_Q\), and total signer weight \(W_Q>3F_Q+2P_Q\).  Hence any two quorums of signer weight
at least \(W_Q-F_Q-P_Q\) in the same committee intersect in correct signer weight; fast quorums of
weight at least \(W_Q-P_Q\) only strengthen this.  Each committee in \(\Committees(\id)\) runs the
single-committee discipline independently, and the election outcome is the deterministic merge of its
local terminal results.  Shared votes are only a wire optimization: per committee they have exactly the
same effect as the corresponding scoped vote.  Before a correct replica final-votes \(x\), its complete
committee-local support, including its first vote and every timeout or ordinary notarization vote, is a
subset of \(\{x\}\).  A correct final voter emits no later first, notarization, or timeout vote in that
same committee instance.

\paragraph{Safety assertions.}
Let \(\mathsf{Correct}\) be the set of correct replicas,
\(\mathsf{Finalized}_j(q)\) the set of non-timeout block ids that replica \(P_j\) has finalized for
slot election \(q\), and \(\closeDecision_j[I].hash=\bot\) when \(P_j\) has not yet decided logical
close instance \(I\).  The following assertions must hold in every reachable state:
\[
\begin{aligned}
\mathsf{ASSERT}_{\mathsf{slot}}\equiv{}&
\forall P_i,P_j\in\mathsf{Correct},\ \forall q=\Slot(a,s,e),\
\forall b\in\mathsf{Finalized}_i(q),\ \forall b'\in\mathsf{Finalized}_j(q):
b=b',\\
\mathsf{ASSERT}_{\mathsf{cut}}\equiv{}&
\forall P_i,P_j\in\mathsf{Correct},\ \forall e,\
\closeDecision_i[\CloseCut(e)].hash\neq\bot\wedge
\closeDecision_j[\CloseCut(e)].hash\neq\bot\\
&\hspace{8em}\Longrightarrow
\closeDecision_i[\CloseCut(e)].hash=\closeDecision_j[\CloseCut(e)].hash,\\
\mathsf{ASSERT}_{\mathsf{record}}\equiv{}&
\forall P_i,P_j\in\mathsf{Correct},\ \forall e,\
\closeDecision_i[\CloseRecord(e)].hash\neq\bot\wedge
\closeDecision_j[\CloseRecord(e)].hash\neq\bot\\
&\hspace{8em}\Longrightarrow
\closeDecision_i[\CloseRecord(e)].hash=\closeDecision_j[\CloseRecord(e)].hash.
\end{aligned}
\]
The first assertion compares finalization, not merely replica-local terminal observations.  The latter
two compare the canonical value hashes decided for the logical close elections across all rounds;
collision resistance then binds equal hashes to equal canonical cuts or frontier maps.

\paragraph{Same election.}
For a fixed \(\Slot(a,s,e)\), first-vote uniqueness, \(\manyVotes(Q,\id)\),
\(\timeoutAllowed(Q,\id)\), and the post-final lock satisfy the committee-local kernel defined above.
By the committee-local certificate-incompatibility lemma, a local fast or final certificate for \(b\)
excludes every
local notarization certificate for \(b'\neq b\), including timeout.  A global fast or final certificate
requires the same value in every election committee, so two conflicting local notarization
certificates exclude every possible present or delayed ordinary global final certificate.  The merge
function finalizes a value only when all committee-local results finalize that same value; one timeout
result dominates any non-final compatible result, and different non-timeout values merge to timeout.
Therefore a committee-local result alone has no global finality or durable release effect.  One
matching merged notarization may finish a slot for drain, but remains unresolved until a close-record
hash is certified.  A correct close voter includes its block on a candidate ledger chain only while its
current validated evidence contains exactly one unresolved notarization and no final or release
evidence.  Known conflicting notarizations resolve to release; a hidden conflicting notarization
remains non-final and is rolled back if another block occupies that logical slot on the certified
ledger.

\paragraph{Release.}
If \(\released_{P_j}(\Slot(a,s,e),b)\), then the certified close state has made \(b\) incapable of
later finalization.  Release follows when the slot is excluded from the certified cut or when the candidate does not occur
on the complete account chains committed by the certified ledger-state root.  A timeout or a known pair
of conflicting notarizations supplies a positive witness permitting such omission.  If a hidden
conflicting notarization is disclosed only after a close record including \(b\) on the certified
ledger has certified, \(b\) remains finalized and the hidden block is released by certified omission.
Certificate-level incompatibility proves that no block from
a timeout- or conflict-released election can later obtain an ordinary global fast or final certificate,
and uniqueness of the certified close hash prevents a competing close finalization.  The
\(\StablePrefix\) lock additionally makes every descendant certificate invalid if its parent chain
contains a released slot, so ancestor finality cannot revive it indirectly.  Bare timeout or conflict
evidence may finish the slot for drain but is not durable release while the close record remains
unsettled.

\paragraph{Adjacent elections.}
\(\Slot(a,s,e)\) and \(\Slot(a,s,e+1)\) share committee \(\committeeAt(e-2)\).  Every global
non-timeout finishing certificate contains the corresponding local certificate from that shared
committee.  Two conflicting local quorums in the shared committee intersect in a correct replica.
That replica retains its slot-vote history across the epoch boundary, and
\(\safeToVote\) forbids its conflicting later vote unless the earlier obligation has been
certifiably released.  A finishing certificate that still preserves the earlier obligation therefore
contradicts the existence of a conflicting adjacent certificate.

\paragraph{Close visibility and cuts.}
A close-cut value becomes available after a report quorum in every committee of
\(\Committees(\CloseCut(e,0))\), and a non-timeout close-cut decision requires a compatible fast
or final result in every close committee.  Every correct close-cut voter verifies a joint report proof and
requires the cut to contain both its current \(visible[e]\) and \(\ReportVisible(e,\rho)\).  Any
pre-close slot election that can affect safety is therefore included either by vote visibility or,
through quorum intersection in each relevant committee, by the signed report proof.  A certified close
cut therefore either records the relevant slot election and waits for it to reach a terminal election
outcome before epoch \(e\) is marked closed, or safely excludes it and provides release after the epoch
closes.  Reaching \(\notarized(b)\) is sufficient for close drain because the close-record phase will
either include \(b\) on a certified account chain or release the slot by omitting it from the certified
ledger.

\paragraph{Safety across close rounds.}
Suppose a valid fast or final certificate for hash \(h\) exists in close round \(I_r\).  Either
certificate decides the logical close instance.  The committee-local certificate-incompatibility lemma
implies that \(I_r\) cannot also have valid timeout evidence or certificate-level conflicting
notarization evidence.  Therefore a correct replica that advances from \(I_r\) without yet learning
the deciding certificate can advance only by carrying \(h\).

In every successor round reached through this live carry, correct replicas first-vote \(h\), even if
another hash is freshly preferred.  Byzantine replicas alone cannot place another value in the
second-look set because their total weight is at most \(F_Q<F_Q+P_Q\), while the second-look threshold
is strictly greater than \(F_Q+P_Q\).  Moreover, once all correct first-vote weight is for \(h\),
\[
\allVotes(Q,I_{r+1})-\maxVotes(Q,I_{r+1})
\le F_Q
\le F_Q+P_Q,
\]
so the ordinary timeout trigger cannot become true solely from Byzantine votes.

Thus no successor round can produce a notarization certificate, timeout certificate, fast certificate,
or final certificate for a value different from \(h\).  By induction, every later close-round fast or
final decision certificate certifies \(h\).

Conversely, if a valid death proof exists for \(I_r\), the committee-local certificate-incompatibility lemma excludes every
fast certificate, and in fact every final certificate, in that round.  Resetting to a fresh hash after
such positive evidence is therefore safe.

The argument does not rely on certificate arrival order.  A valid older-round fast or final certificate
is processed as a decision, even after a replica has entered later rounds.  The instance-wide variable \(\closeDecision[I]\) is
write-once.

\paragraph{Hash-reconstructed close decision.}
For a decided close-cut hash, every correct replica reconstructs the canonical slot set, recomputes the
hash, derives the deciding certificate, and validates the cut through
\(\admissibleHistorical(\CloseCut(e),h_C,C)\).  The certificate attests, under the committee fault
bound, that correct weight freshly validated the report quorum and every start obligation before voting;
those local witness choices need not be serialized or reconstructed after certification.  For a decided
close-record hash, every correct replica reconstructs the canonical frontier map, recomputes the hash
using the preceding close-record hash, and validates the historically admissible decided cut, frozen
close input, slot evidence, complete account chains, ledger transitions, ancestry, and
duplicate-account-slot constraints.  Evidence learned after a vote may change fresh preference but
cannot change the canonical preimage of a decided hash.

\paragraph{Soundness of reconciliation.}
A reconciliation delta has no consensus authority.  It is accepted only relative to a locally retained
canonical base, and only when applying the delta produces a canonical set or frontier map whose hash is
exactly the target and whose applicable fresh, historical, or frozen-input admissibility checks succeed.  A malformed delta therefore either
fails application, produces the wrong hash, or fails semantic validation.  Subject to collision
resistance of \(\hash\), it cannot make a correct replica accept contents different from the contents
committed by the target hash.  Applying a delta never deletes underlying blocks, reports, or votes and
does not alter the replica's current report store, vote store, or derived visibility; it only reconstructs
an immutable historical version indexed by hash.

\paragraph{Disjoint elections.}
For \(e'\ge e+2\), elections need not share a committee.  Every valid vote in epoch \(e'\), however,
is interpreted under the implicit governing context
\(\governingHash(e')=\closeHash[e'-2]\), and the voter must validate the corresponding cumulative
\(\closeState[e'-2]\) and its certified account frontiers.  That state contains only finalized or
released slot outcomes.  If it finalizes block \(b\) as the frontier of account \(a\),
\(\safeFromClose\) permits only blocks of \(a\) whose strict parent chain includes \(b\) and whose
every strict ancestor is finalized.  If the state instead contains \(\ReleasedState(\pi)\), the
certified close decision releases the earlier same-account-slot obligation, while \(\StablePrefix\) forbids a
released block from appearing on a later certified descendant chain.  Since every unresolved
notarization is either finalized or released when its close record certifies, no unresolved obligation
is carried across disjoint committee generations.

A correct replica may vote in epoch \(e'\) only after it has processed the certified close state of
\(e'-2\), including durable \(\closeHash[e'-2]\), \(\ledgerRoot[e'-2]\), \(\closeCut[\cdot]\),
\(\frontier[\cdot,\cdot]\), release evidence, and finalized-slot state.  The close state therefore
participates directly in vote validity rather than only indirectly in committee derivation.

\paragraph{Theorem.}
\(\mathsf{ASSERT}_{\mathsf{slot}}\), \(\mathsf{ASSERT}_{\mathsf{cut}}\), and
\(\mathsf{ASSERT}_{\mathsf{record}}\) hold.  In particular, no two correct replicas finalize
conflicting block ids for the same account slot.  If conflicting slot finality came from the same
election, same-election safety contradicts it; from adjacent elections, adjacent safety contradicts it;
and from elections at distance at least two epochs, disjoint-election safety contradicts it.  Close-cut
and close-record agreement follow from certificate incompatibility, live carry across rounds, and the
write-once logical close decision.


\section{Liveness}

The liveness statement distinguishes a terminal slot-election observation from a durable slot
decision.  A notarization may finish the slot election for drain, but the slot is settled only when an
ordinary fast or final certificate finalizes it or when the certified ledger includes its candidate or
releases the election by omission.  Define
\begin{center}
\begin{adjustbox}{max width=\columnwidth}
\(\displaystyle
\begin{aligned}
\func{settled}(\Slot(a,s,e))\iff{}&
\finalizedSlot[a,s]\neq\bot\\
&\vee\bigl(state[e]=\mathsf{closed}\wedge\closeCut[e]\neq\bot\\
&\qquad\wedge(\Slot(a,s,e)\notin\closeCut[e]\vee
\closeState[e](\Slot(a,s,e))\in\{\FinalizedState(b,\pi),\ReleasedState(\pi)\})\bigr).
\end{aligned}
\)
\end{adjustbox}
\end{center}

\paragraph{Liveness assertion.}
With \(\Diamond\) denoting ``eventually'' over an admissible execution satisfying the liveness
assumptions below, epoch closure is asserted replica-by-replica:
\[
\mathsf{ASSERT}_{\mathsf{epoch\mbox{-}closure}}\equiv
\forall e\ge0,\ \forall P_j\in\mathsf{Correct}:\quad
\Diamond\bigl(state_j[e]=\mathsf{closed}\bigr).
\]
This is a temporal progress assertion, not a finite-prefix invariant: it is evaluated only for
executions with the stated eventual-synchrony, responsiveness, gossip, finite-evidence, retention, and
reconstruction assumptions.
For a logical close instance \(I\in\{\CloseCut(e),\CloseRecord(e)\}\), settled means
\(\closeFinished(I)\); a round-local \(\converged(h)\) observation is not settlement.  Thus
conditional election liveness for RAI is
\begin{center}
\begin{adjustbox}{max width=\columnwidth}
\(\displaystyle
\started[\id]\Rightarrow \Diamond\func{settled}(\id),
\)
\end{adjustbox}
\end{center}
where slot elections use the definition above and logical close instances use
\(\func{settled}(I)\iff\closeFinished(I)\).

\paragraph{Close-liveness assumptions.}
Assume:

\begin{enumerate}
\item after some time the network is synchronous;
\item every correct committee member remains responsive long enough to perform every enabled
      first-vote, second-look, timeout-notarization, and final-vote action;
\item every valid block, report, and vote learned by one correct online replica is eventually made
      available to every correct online replica through duplicate-suppressing reliable gossip;
\item a correct replica that votes for a close hash retains its canonical cut or frontier preimage until
      the close instance decides;
\item decided versions, or valid delta paths from retained checkpoints, remain retrievable from correct
      archival replicas;
\item reconstruction and validation are objective, so every correct replica can eventually reconstruct
      and validate a carried hash;
\item the block, report, and vote universe relevant to elections in the frozen certified cut is finite,
      or eventually stops changing the canonical fresh close record; and
\item fresh construction and canonical serialization are deterministic.
\end{enumerate}

The parameter \(P_Q\) provides quorum-weight slack but does not independently provide crash tolerance for
this liveness theorem.

\paragraph{Common validation of a carry.}
A non-timeout notarization certificate is locally derived from enough signed votes that at least one
correct voter reconstructed and validated the corresponding close version.  Under version
retention, reconciliation availability, and complete atom gossip, every correct replica eventually
reconstructs that same canonical version.  For a carried cut hash, the underlying same-hash committee-local certificates make the historical
admissibility branch permanent; for a carried frontier hash, validation
uses the immutable close input.  Every correct replica can therefore validate and first-vote the carried
hash even if its fresh preference differs.

\paragraph{Relative reconstruction of a carry.}
A correct replica obtains the canonical preimage of a live carried hash by applying an add/remove cut
delta or sparse frontier replacements to a retained base, or by retrieving a complete version when no
usable base is available.  It recomputes the target hash and evaluates fresh admissibility when no
support certificate exists, historical admissibility for a supported cut, or frozen-input admissibility
for a frontier version.  Failure, delay, or corruption of one reconciliation response cannot prevent
progress because the replica may ask another correct holder or obtain the complete canonical version.

\paragraph{Eventual fresh-choice convergence.}
Let \(E_j(t)\) be the relevant valid evidence known to correct replica \(P_j\).  Complete gossip and
the finite-evidence assumption imply that the correct replicas eventually have the same effective
evidence for fresh construction.  Canonical construction therefore gives
\[
\exists T,h^\ast\quad
\forall t\ge T,\ \forall P_j\text{ correct}:
\FreshCloseHash_{P_j}(I)=h^\ast,
\]
whenever no live carry constrains the next first vote.

Correct replicas need not have identical complete sets of admissible hashes.  They need only
eventually validate every live carry and eventually select the same fresh hash in an unlocked round.

For state synchronization and bootstrapping, a joining replica may use the next live close as the latest
target but obtains the durable blocks, reports, and votes and reconstructs the decided cut and frontier
versions for every epoch from genesis onward.  For each epoch it validates all blocks and transitions,
locally derives the deciding close-vote quorums, recomputes the close-cut hash, and recomputes the
close-record hash from the preceding close-record hash and the canonical frontier map.  The terminal
validated close commits transitively to the certified block chains and yields the live committee by
replay.  The archive supplying atoms or reconciliation deltas is not trusted.

A replica starts closing an open epoch only after the previous epoch and its close state are
certified; when it starts closing that epoch, it immediately opens the successor epoch.  Every
responsive participant in epoch \(e\) eventually obtains \(\closeState[e-2]\) before voting.  When
a decided close cut \(C\) is installed, elections in \(C\) are resumed for the drain phase until
they finish, while elections outside \(C\) are not resumed and are settled by certified exclusion.  The complete-gossip and finite-evidence assumptions above govern fresh close-record construction.

\paragraph{Election-core termination.}
Fix an enabled committee instance \((\id,Q)\), and write
\(W=W_Q\), \(F=F_Q\), \(P=P_Q\), \(g=F+P\), and \(q=W-F-P\).
Assume a sufficiently long synchronous period in which every correct member of \(Q\) remains
online long enough to send its first vote, receive every correct first vote, and perform every enabled
second-look or timeout action.  Byzantine members have total weight at most \(F\), and
\(W>3F+2P\).  Then the committee instance eventually derives one local terminal result
\[
L_Q(\id)\in\{\timeoutOutcome,\notarized(x),\fast(x),\final(x)\}.
\]

Let \(F'\le F\) be the actual Byzantine weight.  After all correct first votes have been delivered,
exactly one of the following cases holds.

First, suppose some non-timeout value \(x\) has valid first-vote signer weight strictly greater than
\(g=F+P\).  At most \(F\) of that weight is Byzantine, so the correct first-vote weight supporting
\(x\) is strictly greater than \(P\).  A correct replica emits a non-timeout first vote only after
obtaining the block data or reconstructing the close version and validating \(\admissible(\id,x)\).
Under complete atom gossip, version retention, and reconciliation availability, every correct member of
\(Q\) eventually obtains the data required to validate \(x\).  For a slot value, \(\ExtendsTree\),
\(\safeFromClose\), and the enabled retry guards depend only on the validated block and ancestry,
monotone finalized state, and certified governing close state.  For a close value, admissibility is
stable by definition.  Hence additional valid evidence cannot make this previously validated value
inadmissible, and every correct member eventually performs the committee-local second look and issues a
notarization vote for \(x\).  Correct signer weight is at least \(W-F'\ge W-F\ge W-F-P=q\), so
\(Q\) derives a local notarization certificate for \(x\).

Second, suppose no non-timeout value has first-vote signer weight strictly greater than \(g\).  Once all
correct first votes are present,
\[
\allVotes(Q,\id)\ge W-F'
\qquad\text{and}\qquad
\maxVotes(Q,\id)\le F+P.
\]
Therefore
\[
\allVotes(Q,\id)-\maxVotes(Q,\id)
\ge (W-F')-(F+P)
\ge W-2F-P
>F+P.
\]
Thus \(\timeoutReady(Q,\id)\) becomes true.  Every correct member that has not final-voted in this
committee instance eventually issues the enabled timeout notarization vote; if a correct member has
already final-voted, the corresponding local final result is already obtainable.  Hence in either case
\(Q\) eventually derives a local final result or a local timeout certificate.  This is a
stake-weighted projection of the two cases in Kudzu's Liveness~II lemma
\cite[Lemma~5.10]{shoup2025kudzu}; the additional RAI assumptions above supply stable admissibility and
close-version reconstruction.

Shared first votes count in every applicable committee instance using the signer's frozen weight.
Second-look, timeout, and final votes are nevertheless interpreted per committee: when enabled actions
differ, correct overlap replicas emit scoped votes, and a post-final lock in one committee does not
consume or block the corresponding action in another committee.  Consequently the preceding argument
applies independently to every \(Q\in\Committees(\id)\), even when different committees reach the
second-look threshold for different values.

By the committee-composition lemma, after every committee has a local terminal result, the
deterministic merge immediately finishes the current election round.  For a slot election, notarization
in every committee yields \(\notarized(x)\).  For a close election, matching notarizations yield
\(\converged(x)\), which starts a carried-value successor round rather than certifying the close.
Timeout in one committee and a non-timeout result in another yields timeout; matching final results for
the same value yield \(\final(x)\).  For a close election, that result is a decision.  Different
non-timeout values, including conflicting committee-local final values, merge to timeout and their
signed terminal certificates form conflict death evidence that unlocks a later round.

The committee-local termination argument adapts only Kudzu's two-case vote-counting structure; RAI's
values are derived implicitly as above, and first votes are the first commitment to the derived value.
For close elections, a matching non-timeout notarization certificate creates a live carry, a valid fast
or final certificate decides the logical instance, and timeout or conflict death evidence unlocks a
later round.

\paragraph{Safe close attempts after joint report quorums.}
Before a non-timeout close-cut vote, a correct replica locally selects a joint report proof containing
valid eligible signed reports of combined signer weight at least \(W_Q-F_Q\) in every committee
\(Q\in\Committees(\CloseCut(e,0))\).  It requires its canonical cut version to include
\(\ReportVisible(e,\rho)\), in addition to its current local visibility.  The proof is not transmitted
with the close vote; every receiver derives a suitable proof from its gossiped report store.  A close
attempt need not contain every slot election that may later become visible by gossip.  It must contain
every pre-close slot election that is already safety-obligating, meaning an election whose pre-close
voting evidence could still produce a non-timeout finishing certificate if the epoch were not closed.

Let \(Q\in\Committees(\Slot(a,s,e))\) be any slot committee, and let \(R_Q\subseteq Q\) be the
eligible report signer set used to form a close-cut value in that committee, with
\(\SignerWeight_Q(R_Q)\ge W_Q-F_Q\).  Any non-timeout finishing certificate for
\(\Slot(a,s,e)\) must meet the slot threshold in committee \(Q\).  For the slow threshold, let
\(S_Q\subseteq Q\) be the corresponding signer set, with
\(\SignerWeight_Q(S_Q)\ge W_Q-F_Q-P_Q\).  Then
\[
\SignerWeight_Q(R_Q\cap S_Q)
\ge (W_Q-F_Q)+(W_Q-F_Q-P_Q)-W_Q
=W_Q-2F_Q-P_Q
>F_Q+P_Q.
\]
After removing all Byzantine weight, \(R_Q\cap S_Q\) contains correct signer weight strictly greater
than \(P_Q\), and therefore strictly greater than \(F_Q\) because \(P_Q\ge F_Q\).  Every such correct
signer includes \(\Slot(a,s,e)\) in its report because its
non-timeout vote obligation is unreleased, even if it has already observed a replica-local terminal
outcome.  A fast slot-finishing quorum only increases \(\SignerWeight_Q(S_Q)\), so the same conclusion
holds for fast finalization.  Since the argument holds for every
\(Q\in\Committees(\Slot(a,s,e))\), the joint report proof satisfies
\((a,s)\in\ReportVisible(e,\rho)\).  Admissibility therefore forces every safety-obligating account
slot into the concrete cut before any correct close-cut vote is cast.

Consequently, any close-cut hash stored in the instance-wide close decision safely freezes the cut
\begin{center}
\begin{adjustbox}{max width=\columnwidth}
\(\displaystyle
C=\cutVersion[e,h_C].
\)
\end{adjustbox}
\end{center}
Every safety-obligating pre-close slot election is in \(C\).  A slot election that becomes visible only
after \(h_C\) is certified, but is not in \(C\), is discarded by the certified exclusion and is no
longer an obligation of epoch \(e\).

\paragraph{Visibility and fresh preference across close attempts.}
Vote-derived visibility is monotone, while report-derived visibility is recomputed from the unique
eligible report of each non-equivocating signer.  Because correct replicas gossip all valid reports and
votes, evidence seen by one correct online replica eventually becomes available to every other correct
online replica, and every correct replica eventually excludes the same equivocating report signers.
In an unlocked round, this may change \(\FreshCloseHash_{P_j}(\CloseCut(e))\) or
\(\FreshCloseHash_{P_j}(\CloseRecord(e))\).  It does not delete a retained historical candidate.
Unsupported candidates are freshly revalidated under the current eligible-report state before a new
vote; a cut with persistent certificate support is instead validated by the historical branch and never
loses validity.  A live carry overrides fresh preference whenever its carried version remains admissible.  Under the
finite-evidence and canonical-construction assumptions,
fresh preferences eventually converge whenever no carry constrains the next first vote.

\paragraph{Close-round termination.}
Consider a close round \(I_r\).

If a non-timeout hash \(h\) obtains first-vote signer weight strictly greater than \(F_Q+P_Q\) in an
applicable committee, stable admissibility, atom gossip, and version reconciliation allow every correct
committee member to perform the second look for \(h\).  Thus the popular-candidate case cannot deadlock
merely because another hash has become freshly preferred.

If no non-timeout hash obtains first-vote signer weight strictly greater than \(F_Q+P_Q\), then after
all correct first votes are delivered,
\[
\allVotes(Q,I_r)-\maxVotes(Q,I_r)>F_Q+P_Q,
\]
so the committee-local \(\timeoutReady\) rule becomes enabled.

Therefore every unlocked round eventually:
\begin{itemize}
\item produces a deciding fast or final certificate;
\item produces a live non-timeout carry; or
\item produces a death proof.
\end{itemize}

If it produces a live carry \(h\), every correct replica eventually validates \(h\) and first-votes it
in the successor round.  Since \(P_Q\ge F_Q\),
\[
W_Q-F_Q\ge W_Q-P_Q,
\]
so the correct replicas alone reach the fast threshold and decide \(h\).

If it produces a death proof, fresh voting resumes.  Several unlocked rounds may occur before evidence
convergence, but each terminates.  Eventually an unlocked round begins after all correct replicas
compute the same fresh hash \(h^\ast\); the correct first-vote weight then reaches the fast threshold
and decides \(h^\ast\).

\paragraph{Draining the decided cut.}
Once \(\closeDecision[\CloseCut(e)]=(r,h_C)\), replicas reconstruct and install
\begin{center}
\begin{adjustbox}{max width=\columnwidth}
\(\displaystyle
C=\cutVersion[e,h_C],\qquad \closeCut[e]=C.
\)
\end{adjustbox}
\end{center}
For every \(q=\Slot(a,s,e)\in C\), the historically admissible certified cut preserves the slot's start
obligation, and installation initializes \(T_{\mathrm{election}}[q]\) if necessary.  Thus either the slot election has already finished, or it
is resumed with a running timeout during the close-drain phase.  By election-core termination, each
such included slot election eventually
finishes, possibly with \(\notarized(b)\) or \(\timeoutOutcome\).  For every \(\Slot(a,s,e)\notin C\), the election is not
resumed after the certified cut and is settled by exclusion once epoch \(e\) closes.

When epoch \(e\) enters closing, \(\StartClosing(e)\) immediately opens epoch \(e+1\), subject to
the invariant that epoch \(e-1\) is already closed.  The decided cut then fixes
\(\mathsf{CloseInput}_e=(\closeHash[e-1],\frontier[e-1,\cdot],C)\).  After all elections in \(C\) have
finished, each replica applies the deterministic resolver only to those elections and performs ancestry
closure by close-local replay from \(\closeState[e-1]\), deriving a candidate frontier for every
account.  Ordinary finalizations from elections outside \(C\), including epoch-\(e+1\) elections, do not
change the epoch-\(e\) candidate frontier.  A global fast or final certificate requires its block
and every validated ancestor on its parent chain to occur on the candidate ledger only when that chain
has a certified-stable prefix; locally derived timeout or conflicting-notarization evidence omits the
slot; and exactly one currently known non-timeout notarization, with no known final or release evidence,
may place its block on the candidate account chain.  There is no tie-breaking when a conflict is known.

The evidence known to a correct replica grows monotonically.  If no live carry constrains the next
round, complete atom gossip and finite evidence eventually give all correct online replicas the same
effective evidence for fresh construction, so they derive the same canonical frontier map \(\Phi_e\)
and the same close-record hash
\[
\hash(\mathsf{RAI/CloseRecord}\parallel e\parallel\closeHash[e-1]
\parallel\operatorname{canon}(\Phi_e)).
\]
If a close-record fast or final decision is derived in any round, the instance-wide decision is authoritative:
every block on the decided account chains finalizes and omitted epoch-\(e\) elections are released, even
if conflicting notarization or timeout votes are learned later.  Once
\(\closeDecision[\CloseRecord(e)]=(r,h_F)\), the epoch loop reconstructs the decided frontier version,
validates and installs \(\closeState[e]\) and \(\frontier[e,\cdot]\), sets
\(\ledgerRoot[e]\gets\hash(\operatorname{canon}(\Phi_e))\) and \(\closeHash[e]\gets h_F\), derives
\(\committee[e]\) from the installed balances and representatives, and calls \(\FinishClosing(e)\).  Thus epoch \(e+1\) opens when epoch \(e\)
starts closing, and epoch \(e\) eventually becomes closed after its certified cut has reached terminal
outcomes and its compact close record has been decided.


\paragraph{Theorem.}
Under the assumptions above, every started slot election is eventually settled and every started
logical close instance eventually decides.  A close round either decides, creates a live carry, or
produces a death proof.  A live carry is validated and fast-decided in its successor round; a death
proof unlocks canonical fresh voting, which eventually converges and reaches the fast threshold.
Valid fast or final certificates are processed as decisions in every round, including after a replica
has entered a later round.  The decided close cut is
drained, the hash-reconstructed decided close record installs the certified account frontiers and all
derived finalized and released states, and the successor epochs continue to open and close during
periods of synchrony.  Consequently,
\(\mathsf{ASSERT}_{\mathsf{epoch\mbox{-}closure}}\) holds.

\balance
\begin{thebibliography}{1}

\bibitem{shoup2025kudzu}
Victor Shoup, Jakub Sliwinski, and Yann Vonlanthen.
\newblock 2025.
\newblock Kudzu: Fast and Simple High-Throughput BFT.
\newblock \emph{arXiv preprint arXiv:2505.08771v2} (29 September 2025).

\end{thebibliography}
\end{document}
